\documentclass[12pt]{article}
\title{Digest of The New York Times}
\author{The New York Times}

% % \definecolor{dark-green}{rgb}{0,0.25,0}
% \usepackage[utf8]{inputenc}

\usepackage{fontspec,xunicode,xltxtra}
\setmainfont[Ligatures=Required,Mapping=tex-text]{Georgia}
% \setmainfont[Ligatures=Required,Mapping=tex-text]{Palatino Linotype}
\setsansfont[Ligatures=Required,Mapping=tex-text]{Myriad Pro}
\setmonofont[Ligatures=Required,Mapping=tex-text]{Courier New}
\usepackage{xeCJK}
\setCJKmainfont{宋体}
\setCJKmonofont{黑体}

% \usepackage{indentfirst}


\usepackage[dvips,dvipsnames,svgnames]{xcolor}
\definecolor{light-gray}{gray}{0.95}

%生成PDF的链接
\usepackage[dvipdfmx]{hyperref}
\hypersetup{
    bookmarks=true,         % show bookmarks bar?
    bookmarksopen=true,
    pdfpagemode=UseNone,    % options: UseNode, UseThumbs, UseOutlines, FullScreen
    pdfstartview=FitH,
    pdfborder=1,
    pdfhighlight=/P,
    pdfauthor={wuxch},
    unicode=true,           % non-Latin characters in Acrobat’s bookmarks
    colorlinks,             % false: boxed links; true: colored links
    linkcolor=blue,         % color of internal links
    citecolor=blue,        % color of links to bibliography
    filecolor=magenta,      % color of file links
    urlcolor=cyan           % color of external links
}
% 使用目录
\usepackage{makeidx}
% 使用图形
\usepackage{graphicx}


\usepackage{titletoc}
\dottedcontents{section}[1.5em]{}{2.3em}{1pc}

% % 页眉页脚的设定
% \usepackage{fancyhdr}
% \pagestyle{fancy}
% \lhead{\includegraphics[height=0.2in]{zte_logo}}

% \usepackage{footmisc}
\usepackage{pfnote}
\usepackage{fnpos}

% % 嵌入的代码显示
% \usepackage{listings}
% \lstset{language=C++, breaklines, extendedchars=false}
% \lstset{keywordstyle=\color{blue},
%         commentstyle=\color{SeaGreen},
%         stringstyle=\ttfamily,
%         showstringspaces=false,
%         backgroundcolor=\color{light-gray}}
% \lstset{basicstyle=\small}

\usepackage{booktabs}

\usepackage{tipa}
\usepackage{phonetic}
\usepackage[gen]{eurosym}

% \usepackage{CJK}
\usepackage{makeidx}
% \usepackage{glossary}

\usepackage{fancyhdr}

\usepackage{lastpage}
\cfoot{\textsf{\thepage\ of \pageref*{LastPage}}}

\pagestyle{fancy}
\fancyhead{}
\renewcommand{\headrulewidth}{0pt}
\pagenumbering{arabic}

\usepackage{graphicx}
\usepackage{picinpar}

\setlength{\parindent}{0in}
\setlength{\parskip}{2.5ex plus4mm minus3mm}
% \setlength{\parskip}{.3\baselineskip}

\setcounter{section}{0}

% bibtex
\usepackage[super]{natbib}
% \usepackage[globalcitecopy,labelstoglobalaux]{bibunits}
% \makeatletter
% \renewcommand{\@bibunitname}{\jobname.\the\@bibunitauxcnt}
% \makeatother

% \renewcommand{\refname}{缩略语}
% \renewcommand{\refname}{}

\makeatletter                   % Reference list option change
\renewcommand\@biblabel[1]{#1:  } % from [1] to 1: 
\makeatother                    %

\newcommand{\citenumfont}[1]{\textsf{#1}}

% \usepackage{tipa}
% \usepackage{phonetic}

\newcommand{\phonetic}[1]{{\fontspec[Colour=CC00CC]{Microsoft Sans Serif} #1}}


%%% Local Variables:
%%% TeX-master: "medical_trivia"
%%% End:

\usepackage{config}

\makeindex
\begin{document}
\date{}
\thispagestyle{empty}
\renewcommand\contentsname{\textsf{Digest of The New York Times}}
{\footnotesize\textsf{\tableofcontents}}
\clearpage
\setcounter{page}{1}

\section{Another Peril\cite{peril} in War Zones: Sexual Abuse by Fellow G.I.'s}

\lettrine{C}{apt}\mycalendar{Dec.'09}{28}. Margaret H.~White began a relationship with a warrant
officer while both were training to be deployed to Iraq. By the time they arrived this year at Camp
Taji, north of here, she felt what she called ``creepy vibes'' and tried to break it off.

In the claustrophobic confines of a combat post, it was not easy to do. He left notes on the door to
her quarters, alternately pleading and menacing. He forced her to have sex, she said. He asked her
to marry him, though he was already married. He waited for her outside the women's latrines or her
quarters, once for three hours.

``It got to the point that I felt safer outside the wire,'' Captain White said, referring to
operations that take soldiers off their heavily fortified bases, ``than I did taking a shower.''

Her ordeal ended with the military equivalent of a restraining order and charges of stalking against
the officer. It is one case that highlights the new and often messy reality the military has had to
face as men and women serve side by side in combat zones more than ever before.

Sexual harassment and sexual assault, which the military now defines broadly to include not only
rape but also crimes like groping and stalking, continue to afflict the ranks, and by some measures
are rising. While tens of thousands of women have served in Iraq and Afghanistan, often in combat,
often with distinction, the integration of men and women in places like Camp Taji has forced to the
surface issues that commanders rarely, if ever confronted before.

The military -- belatedly, critics say -- has radically changed the way it handles sexual abuse in
particular, expanding access to treatment and toughening rules for prosecution. In the hardships of
war, though, the effects of the changes remain unclear.

The strains of combat, close quarters in remote locations, tension and even boredom can create the
conditions for abuse, even as they hinder medical care for victims and legal proceedings against
those who attack them.

Captain White said she had feared coming forward, despite having become increasingly despondent and
suffered panic attacks, because she was wary of she-said-he-said recriminations that would
reverberate through the tightknit military world and disrupt the mission. Despite the military's
stated ``zero tolerance'' for abuse or harassment, she had no confidence her case would be taken
seriously and so tried to cope on her own, Captain White said.

A Pentagon-appointed task force, in a report released this month, pointedly criticized the
military's efforts to prevent sexual abuse, citing the ``unique stresses'' of deployments in places
like Camp Taji. ``Some military personnel indicated that predators may believe they will not be held
accountable for their misconduct during deployment because commanders' focus on the mission
overshadows other concerns,'' the report said.

That, among other reasons, is why sexual assault and harassment go unreported far more often than
not. ``You're in the middle of a war zone,'' Captain White said, reflecting a fear many military
women describe of being seen, somehow, as harming the mission.

``So it's kind of like that one little thing is nothing compared with 'There is an I.E.D. that went
off in this convoy today and three people were injured,' '' she said, referring to an improvised
explosive device.

Common Fears

By the Pentagon's own estimate, as few as 10 percent of sexual assaults are reported, far lower than
the percentage reported in the civilian world. Specialist Erica A.~Beck, a mechanic and gunner who
served in Diyala Province in Iraq this summer, recalled a sexual proposition she called
``inappropriate'' during her first tour in the country in 2006-7. ``Not necessarily being vulgar,
but he, you know, was asking for favors,'' she said.

She did not report it, she said, because she feared that her commanders would have reacted harshly
-- toward her.

``It was harassment,'' she said. ``And because it was a warrant officer, I didn't say anything. I
was just a private.''

Her fears were common, according to soldiers and advocates who remain skeptical of the military's
efforts to address abuse. A report last year by the Government Accountability Office concluded that
victims were reluctant to report attacks ``for a variety of reasons, including the belief that
nothing would be done or that reporting an incident would negatively impact their careers.''

When Sgt. Tracey R.~Phillips told a superior about an unwanted sexual advance from a private the
night their unit arrived in Iraq in May, the accusations unleashed a flurry of charges and
countercharges, an initial investigation of her on charges of adultery, a crime in the military
justice system, and, according to her account, violations by her commanders of the new procedures
meant to ease reporting of abuse.

In the end, she was kicked out of Iraq and the Army itself, while the private remained on duty here.

The military disputed her account but declined to state the reasons for sending her out of Iraq. Her
paperwork showed that she received an honorable discharge, though with ``serious misconduct'' cited
as the reason. The so-called misconduct, she said, stemmed from the Army's allegation that she had
had an inappropriate relationship with the private she accused. She denied that.

``If I would have never, ever, ever said anything, I wouldn't be sitting here,'' she said in an
interview at her parents' home near San Antonio. ``I'd still be in Iraq.''

At bases around Iraq, many said that acceptance and respect for women in uniform were now more
common than the opposite. In part, they said, that reflects a sweeping change in military culture
that has accompanied the rise of women through the ranks and into more positions once reserved for
men.

``It's not tolerated -- it's just not,'' said Lt.~Brenda L.~Beegle, a married military police
officer, referring to sexual harassment and abuse.

In an interview at Liberty Base, near Baghdad's airport, she said: ``Everyone has heard stories
about bad things that have happened. I've never had an issue.''

Although exact comparisons to the civilian world are difficult because of different methods of
defining and reporting abuse, Pentagon officials and some experts say that the incidence of abuse in
the military appears to be no higher than in society generally, and might be lower. It appears to be
even lower in combat operations than at bases in the United States, because of stricter discipline
and scrutiny during deployments, as well as restrictions on alcohol, which is often a factor in
assaults, for example, on college campuses.

Complaints Increase

The number of complaints, though, is rising. Across the military, there were 2,908 reported cases of
sexual abuse involving service members as victims or assailants, in the fiscal year that ended in
September 2008, the last year for which the Pentagon made numbers available. That was an 8 percent
increase from the previous year, when there were 2,688.

In the turbulent regions from Egypt to Afghanistan where most American combat troops are now
deployed, the increase in reported cases was even sharper: 251 cases, compared with 174 the year
before, a 44 percent increase. The number in Iraq rose to 143, from 112 the year before. Everyone
agrees that those represent only a fraction of the instances of assault, let alone harassment.

``A woman in the military is more likely to be raped by a fellow soldier than killed by enemy fire
in Iraq,'' Representative Jane Harman, a Democrat from California, said at a Congressional hearing
this year, repeating an assertion she has made a refrain in a campaign of hers to force the military
to do more to address abuses.

At least 10 percent of the victims in the last year were men, a reality that the Pentagon's task
force said the armed services had done practically nothing to address in terms of counseling,
treatment and prosecution. Men are considered even less likely to report attacks, officials said,
because of the stigma, and fears that their own sexual orientation would be questioned. In the
majority of the reported cases, the attacker was male.

Senior Pentagon officials argued that the increase in reports did not necessarily signify a higher
number of attacks. Rather, they said, there is now a greater awareness as well as an improved
command climate, encouraging more victims to come forward.

``We believe the increase in the number of reported cases means the department is capturing a
greater proportion of the cases that occurred during the year, which is good news,'' said the
Pentagon's senior official overseeing abuse policies, Kaye Whitley.

The military can no more eradicate sexual abuse than can society in general, but soldiers, officers
and experts acknowledge that it is particularly harmful when soldiers are in combat zones, affecting
not only the victims but also, as the military relies more than ever on women when the nation goes
to war, the mission.

``For the military the potential costs are even higher as it can also negatively impact mission
readiness,'' the Pentagon's annual report on sexual abuse said, referring to sexual violence.
``Service members risk their lives for one another and bear the responsibility of keeping fellow
service members out of harm's way. Sexual assault in the military breaks this bond.''

Even investigations into accusations, which are often difficult to prove, can disrupt operations. In
Sergeant Phillips's case, she was relieved of her duties leading a squad of soldiers refueling
emergency rescue helicopters and other aircraft at Camp Kalsu, south of Baghdad.

Cases like hers suggest that the vagaries of sex and sexual abuse, especially in combat zones,
continue to vex commanders on the ground, despite the transformation of the military's policies.

The majority of sexual abuse allegations end with no prosecution at all. Of 2,171 suspects of
investigations that were completed during the fiscal year that ended in September 2008, only 317
faced a court-martial. Another 515 faced administrative punishments or discharges. Nearly half of
the completed investigations lacked evidence or were ``unsubstantiated or unfounded.''

The Pentagon, facing criticism, maintains that it has transformed the way it handles sexual abuse.
In the wake of the invasions of Afghanistan and Iraq, as well as highly publicized cases and
revelations of rampant abuse at the Air Force Academy in 2003, the Pentagon created a single agency
to oversee the issue and rewrote the rules of reporting, treatment and prosecution. Beginning in
October 2007, the Uniform Code of Military Justice expanded the provision that once covered rape --
Article 120 -- to include other offenses, like indecent exposure and stalking.

The Army, which has provided the bulk of the forces in Iraq, has increased the number of
investigators and lawyers trained to investigate accusations. Most bases now have kits to collect
forensic evidence in rape cases, which was not the case immediately after the invasion in 2003.

Larger field hospitals in Balad and Mosul now have the same type of sexual assault nurse examiners
widely used in the civilian world, as well as a dozen other examiners who are not nurses but are
trained to conduct forensic examinations.

The military has set up a system of confidential advisers women can turn to who are outside the
usual chain of command -- an avenue Sergeant Phillips said she had been denied.

If they want to, the women can now seek medical treatment and counseling without setting off a
criminal investigation. And all the services have started educational programs to address aspects of
a hierarchical warrior culture that some say contributes to hostility toward women. Posters for the
campaign blanket bulletin boards in offices, chow halls and recreational buildings on bases across
Iraq.

The military's efforts, however well intentioned, are often undermined by commanders who are
skeptical or even conflicted, suspicious of accusations and fearful that reports of abuse reflect
badly on their commands. The Pentagon task force also reported that victims of assault did not come
forward because they might ``have engaged in misconduct for which they could be disciplined, such as
under-age drinking, fraternization or adultery.''

Marti Ribeiro, then an Air Force sergeant, said she was raped by another soldier after she stepped
away from a guard post in Afghanistan in 2006 to smoke a cigarette, a story first recounted in ``The
Lonely Soldier,'' a book by Helen Benedictabout women who served in Iraq and elsewhere. When she
went to the abuse coordinator, she was threatened with prosecution for having left her weapon and
her post.

``I didn't get any help at all, let alone compassion,'' said Ms.~Ribeiro, who has since retired and
joined the Service Women's Action Network, a new advocacy organization devoted to shaping the
Pentagon's policy.

The hardships of combat operations often compound the anguish of victims and complicate
investigations, as well as counseling and treatment. The Government Accountability Office suggested
that the ``unique living and social circumstances'' of combat posts heightened the risk for assault.
Both the G.A.O. and the Pentagon's task force found that, despite the Pentagon's policy, remote
bases did not have adequate medical and mental health services for victims. The task force also
found that abuse coordinators and victim advocates were often ill trained or absent.

As a result, victims often suffer the consequences alone, working in the heat and dust, living in
trailers surrounded by gravel and concrete blast walls, with nowhere private to retreat to. In
Captain White's case, she had to work and live beside the man who assaulted and stalked her until
their deployment ended in August and they both went home.

``You're in such a fishbowl,'' she said. ``You can't really get away from someone. You see him in
the chow hall. You see him in the gym.''

The Danger Nearby

Captain White's case is typical of many here, according to military lawyers and experts, in that she
knew the man she said assaulted her, circumstances that complicated the investigation and
prosecution.

She had dated the warrant officer when they arrived in Fort Dix, N.J., for predeployment training
with the 56th Stryker Combat Team. The newly revised article of the Uniform Code of Military Justice
says that ``a current or previous dating relationship by itself'' does not constitute consent.

Once at Camp Taji, a sprawling base just north of Baghdad, she grew troubled by his behavior. He
cajoled her with presents and sent her e-mail messages. She said that for fear of running into him,
she stopped drinking water after 7 p.m. so she would not have to go to the latrine at night alone.

She never came forward herself. Her case came to light only when military prosecutors questioned her
about another investigation involving the warrant officer. He was ultimately charged with 19
offenses, said Lt.~Col.~Philip J.~Smith, a spokesman for the division that oversaw operations in
central Iraq. The charges included seven counts of fraternization and two of adultery, interfering
with an investigation and, in Captain White's case, stalking.

After their deployment ended in September, the officer pleaded guilty and resigned from the Army in
lieu of prosecution, Colonel Smith said.

Captain White said that she was satisfied with the legal outcome of her case, though her account of
it highlighted the emotional strains that sexual abuse causes.

``I'm not saying that I handled it the best way,'' she said in an interview after her own retirement
from the Army, ``but I handled it at the time and in the situation what I thought was the best way,
which was just to keep my head down, keep going -- which was kind of an Army thing to say: Drive
on.''

\section{Elite U.S.~Force Expanding Hunt in Afghanistan}

\lettrine{S}{ecretive}\mycalendar{Dec.'09}{28} branches of the military's Special Operations forces
have increased counterterrorism missions against some of the most lethal groups in Afghanistan and,
because of their success, plan an even bigger expansion next year, according to American commanders.

The commandos, from the Army's Delta Force and the Navy's classified Seals units, have had success
weakening the network of Sirajuddin Haqqani, the strongest Taliban warrior in eastern Afghanistan,
the officers said. Mr.~Haqqani's group has used its bases in neighboring Pakistan to carry out
deadly strikes in and around Kabul, the Afghan capital.

Guided by intercepted cellphone communications, the American commandos have also killed some
important Taliban operatives in Marja, the most fearsome Taliban stronghold in Helmand Province in
the south, the officers said. Marine commanders say they believe that there are some 1,000 fighters
holed up in the town.

Although President Obama and his top aides have not publicly discussed these highly classified
missions as part of the administration's revamped strategy for Afghanistan and Pakistan, the
counterterrorism operations are expected to increase, along with the deployment of 30,000 more
American forces in the next year.

The increased counterterrorism operations over the past three or four months reflect growth in every
part of the Afghanistan campaign, including conventional forces securing the population, other
troops training and partnering with Afghan security forces, and more civilians to complement and
capitalize on security gains.

American commanders in Afghanistan rely on the commando units to carry out some of the most
complicated operations against militant leaders, and the missions are never publicly acknowledged.

The commandos are the same elite forces that have been pursuing Osama bin Laden, captured Saddam
Hussein in Iraq in 2003 and led the hunt that ended in 2006 in the death of Abu Musab al-Zarqawi,
the leader in Iraq of the insurgent group Al Qaeda in Mesopotamia.

In recent interviews here, commanders explained that the special-mission units from the Joint
Special Operations Command were playing a pivotal role in hurting some of the toughest militant
groups, and buying some time before American reinforcements arrived and more Afghan security forces
could be trained.

``They are extremely effective in the areas where we are focused,'' said one American general in
Afghanistan about the commandos, speaking on the condition of anonymity because of the classified
status of the missions.

Gen.~David H.~Petraeus, who is in charge of the military's Central Command, mentioned the increased
focus on counterterrorism operations in testimony before the Senate Foreign Relations Committee on
Dec.~9. But he spoke more obliquely about the teams actually conducting attacks against hard-core
Taliban extremists, particularly those in rural areas outside the reach of population centers that
conventional forces will focus on.

``We actually will be increasing our counterterrorist component of the overall strategy,'' General
Petraeus told lawmakers. ``There's no question you've got to kill or capture those bad guys that are
not reconcilable. And we are intending to do that, and we will have additional national mission
force elements to do that when the spring rolls around.''

Senior military officials say it is not surprising that the commandos are playing such an important
role in the fight, particularly because Gen.~Stanley A.~McChrystal, the senior American and NATO
officer in Afghanistan, led the Joint Special Operations Command for five years.

In addition to the classified American commando missions, military officials say that other NATO
special operations forces have teamed up with Afghan counterparts to attack Taliban bomb-making
networks and other militant cells.

About six weeks ago, allied and Afghan special operations forces killed about 150 Taliban fighters
in several villages near Kunduz, in northern Afghanistan, a senior NATO military official said.

Some missions have killed Taliban fighters while searching for Pfc. Bowe R.~Bergdahl, who was
reported missing on June 30 in eastern Afghanistan. The Taliban in July posted a video on jihadist
Web sites in which the soldier identified himself and said that he had been captured when he lagged
behind on a patrol. A second video was released on Friday.

``We've been hitting them hard, but I want to be careful not to overstate our progress,'' said the
NATO official, speaking on the condition of anonymity in order to describe the operations in detail.
``It has not yet been decisive.''

In Helmand, more than 10,000 Marines, as well as Afghan and British forces, are gearing up for a
major confrontation in Marja early next year. Brig. Gen.~Larry Nicholson, the senior Marine
commander in the south, said in a recent interview, ``The overt message we're putting out is, Marja
is next.''

General Nicholson said there were both ``kinetic and nonkinetic shaping operations'' under way. In
military parlance that means covert operations, including stealthy commando raids against specific
targets, as well as an overt propaganda campaign intended to persuade some Taliban fighters to
defect.

Military officials say the commandos are mindful of General McChrystal's directive earlier this year
to take additional steps to prevent civilian casualties.

In February, before General McChrystal was named to his current position, the head of the Joint
Special Operations Command, Vice Adm. William H.~McRaven, ordered a halt to most commando missions
in Afghanistan, reflecting a growing concern that civilian deaths caused by American firepower were
jeopardizing broader goals there.

The halt, which lasted about two weeks, came after a series of nighttime raids by Special Operations
troops killed women and children, and after months of mounting outrage in Afghanistan about
civilians killed in air and ground attacks. The order covered all commando missions except those
against the top leaders of the Taliban and Al Qaeda, military officials said.

Across the border in Pakistan, where American commandos are not permitted to operate, the Central
Intelligence Agency has stepped up its missile strikes by Predator and Reaper drones on groups like
the Haqqani network.

But an official with Pakistan's main spy agency, the Inter-Services Intelligence directorate, or
I.S.I., said there had also been more than 60 joint operations involving the I.S.I. and the C.I.A.
in the Federally Administered Tribal Areas and Baluchistan in the past year.

The official said the missions included ``snatch and grabs'' -- the abduction of important militants
-- as well as efforts to kill leaders. These operations were based on intelligence provided by
either the United States or Pakistan to be used against the Taliban and Al Qaeda, the official said.

``We can expect to see more U.S.~action against Haqqani,'' a senior American diplomat in Pakistan
said in a recent interview.

The increasing tempo of commando operations in Afghanistan has caused some strains with other
American commanders. Many of the top Special Operations forces, as well as intelligence analysts and
surveillance aircraft, are being moved to Afghanistan from Iraq, as the Iraq war begins to wind
down.

``It's caused some tensions over resources,'' said Lt.~Gen.~Charles H.~Jacoby Jr., the
second-ranking commander in Iraq.

\section{A Record Budget Stirs Debt Worries in Japan}

\lettrine{T}{he}\mycalendar{Dec.'09}{28} Japanese government approved on Friday a record
trillion-dollar budget for the next fiscal year that encompasses ambitious welfare outlays to help
households cope with the country's deep economic woes, but the scale of new spending could renew
investor jitters about the government's burgeoning debt.

Prime Minister Yukio Hatoyama, who took office in September after landmark elections that ended an
era of single-party rule in Japan, has been trying to strike a balance between a budget big enough
to kick-start economic recovery while keeping new lending in check.

He has been distracted, however, by mishandled negotiations over the fate of a United States
military air base in Japan and by a campaign finance scandal. On Thursday, Mr.~Hatoyama apologized
to the nation after two aides were charged with falsely reporting donations, though he denied any
wrongdoing on his part.

Addressing the country for the second night in a row on Friday, Mr.~Hatoyama, who leads the
Democratic Party, urged the public to stay focused on the task of rebuilding an economy recovering
from its worst recession since World War II.

The record \$1 trillion budget for the year starting in April reorients spending to households by
allocating more to welfare and education at the expense of public works projects, where the ousted
Liberal Democrats had poured a large amount of government funds.

Mr.~Hatoyama hopes the generous welfare spending will encourage households to consume more, offering
a much-needed boost to the economy.

``Together with all of you, I want to build a better Japan, a new Japan,'' Mr.~Hatoyama said at a
news conference. ``I have adhered to the principle that people matter more than concrete,'' he said.

Japan will issue fresh debt worth a record \$485 billion, in line with an earlier estimate, he said.
The new borrowing brings Japan's public debt to about \$9.4 trillion, or 181 percent of gross
domestic product, at the end of March 2011, by far the highest in the industrialized world.

Dogging Mr.~Hatoyama's budgeting efforts has been a plunge in tax revenue as Japan's export-oriented
economy struggles with a falloff in international trade brought on by the global economic crisis.

Tax receipts, at about \$405 billion, are expected to make up less than half the government's
budget, forcing the government to borrow more than it receives in revenues -- another postwar first.

Citing tough finances, the government abandoned a campaign pledge to cut an unpopular tax on
gasoline, and it has backtracked on promises to abolish tolls on Japanese highways. But the
government stuck with its ambitious social agenda, which includes cash payments to child-rearing
households and free public high school education.

Though most voters had appeared to be willing to give Mr.~Hatoyama time to deliver on his promises,
public support for him is waning. A public opinion poll recently showed his approval rating had
fallen below 50 percent, from a post-election high of 71 percent.

The campaign financing scandal has further eroded his credibility. At first, the allegations of
improper accounting raised few anxieties among voters because much of the money involved was from
Mr.~Hatoyama's own coffers or from his mother, an heiress to the fortune of the Japanese tire maker
Bridgestone.

But the charges have increasingly painted Mr.~Hatoyama as a leader out of touch with the economic
plight of average Japanese at a time of high unemployment and persistent deflation.

The falling support puts the Democrats in a potentially precarious position ahead of elections for
Parliament's Upper House, which must be held next year. The Democrats hope to clinch a majority in
that chamber to match their control of Parliament's powerful Lower House.

\section{Passengers' Quick Action Halted Attack}

\lettrine{D}{espite}\mycalendar{Dec.'09}{28} the billions spent since 2001 on intelligence and
counterterrorism programs, sophisticated airport scanners and elaborate watch lists, it was
something simpler that averted disaster on a Christmas Day flight to Detroit: alert and courageous
passengers and crew members.

During 19 hours of travel, aboard two flights across three continents, law enforcement officials
said, Umar Farouk Abdulmutallab bided his time. Then, just as Northwest Flight 253 finally began its
final approach to Detroit around noon on Friday, he tried to ignite the incendiary powder mixture he
had taped to his leg, they said.

There were popping sounds, smoke and a commotion as passengers cried out in alarm and tried to see
what was happening. One woman shouted, ``What are you doing?'' and another called out, ``Fire!''

And then history repeated itself. Just as occurred before Christmas in 2001, when Richard C.~Reid
tried to ignite plastic explosives hidden in his shoe on a trans-Atlantic flight, fellow passengers
jumped on Mr.~Abdulmutallab, restraining the 23-year-old Nigerian.

Jasper Schuringa, a Dutch film director seated in the same row as Mr.~Abdulmutallab but on the other
side of the aircraft, saw what looked like an object on fire in the suspect's lap and ``freaked,''
he told CNN.

``Without any hesitation, I just jumped over all the seats,'' Mr.~Schuringa said, in an account that
other passengers confirmed.``I was thinking, Oh, he's trying to blow up the plane. I was trying to
search his body for any explosive. I took some kind of object that was already melting and smoking,
and I tried to put out the fire and when I did that I was also restraining the suspect.''

Mr.~Schuringa said he had burned his hands slightly as he grappled with Mr.~Abdulmutallab, aided by
other passengers among the 289 on board, and began to shout for water.

``But then the fire was getting worse, so I grabbed the suspect out of the seat,'' Mr.~Schuringa
said. Flight attendants ran up with fire extinguishers, doused the flames and helped Mr.~Schuringa
walk Mr.~Abdulmutallab to first class, where he was stripped, searched and locked in handcuffs.

``The whole plane was screaming -- but the suspect, he didn't say a word,'' Mr.~Schuringa said.

He shrugged off praise for his swift action, which he said was reflexive. ``When you hear a pop on
the plane, you're awake, trust me,'' he said. ``I just jumped. I didn't think. I went over there and
tried to save the plane.''

In an affidavit filed in court, an F.B.I. agent said that Mr.~Abdulmutallab stayed in the bathroom
for 20 minutes before the attempt, returned to his seat, told his seatmates that his stomach was
upset and covered himself with a blanket. It was then that the smoke and popping sounds began.

After he was subdued and the fire extinguished, a flight attendant asked him what had been in his
pocket, and he answered, ``explosive device,'' the affidavit said. The powder was identified by the
F.B.I. as PETN, a high explosive.

The close call was followed by several tense hours as counterterrorism officials checked on other
United States-bound flights to determine whether more planes were targets, as in the thwarted 2006
plot to smuggle liquid explosives aboard multiple flights leaving from Britain.

They found no immediate signs that other flights were in danger, officials said. They tightened
airport security, ordering new restrictions on carry-on luggage and passenger movement inside the
cabin, but did not elevate the nation's threat level, which has been at orange since 2006.

Dozens of investigators led by the Federal Bureau of Investigation were working Saturday to
understand exactly how a passenger managed to get PETN and a syringe of chemicals aboard the flight.
Intelligence agencies were studying intercepted communications to see whether clues were missed and
to assess whether the incident could presage more attacks.

David Schilke, 49, of Livonia, Mich., who works in the information technology department at the Ford
Motor Company, was traveling home from Moscow with his wife, Iliana, and their 5-year-old son,
sitting two rows behind the suspect. He said he heard a pop, and then someone asking for water and
screams coming from the rows in front of him. The fire, he said, lasted for a full minute.

``The guy wasn't fighting or doing anything,'' Mr.~Schilke said. ``He was just sitting there in the
flames. I was shocked that he would do that.'' He added that he was surprised at how little panic
there was. Many passengers who were farther away thought the pops were from fireworks, he said.

Richard Griffith, 41, of Pontiac, Mich., who said he had been sitting in the back of the plane
during the episode, praised the crew for its professionalism in preventing panic.

Mr.~Griffith said the passenger who had been sitting next to the suspect told him the suspect got up
once midflight to use the bathroom and returned to the bathroom about 20 or 30 minutes before the
attempt, apparently to brush his teeth. Otherwise, he said, ``He just sat there; he didn't talk to
nobody.''

The episode, which riveted the attention of President Obama on vacation in Hawaii and prompted
counterterrorism officials to rush back to work, capped a year in which plots of violence inside the
United States have surged. The attempt appeared to underscore the continuing determination of Muslim
militants to kill Americans more than eight years after the Sept.~11 attacks.

Passengers transferring from foreign flights at the Amsterdam airport, where Mr.~Abdulmutallab
changed planes and boarded the flight bound for Detroit, are required to be screened by security
there before taking off on another flight, an airport spokeswoman said Saturday. She could not
confirm the details in Mr.~Abdulmutallab's case but said he was presumably subject to that sort of
screening.

Investigators planned to interview all the passengers on the suspect's flights and to look over any
security-camera video footage of him, a law enforcement official said.

Mr.~Abdulmutallab apparently left Murtala Muhammed Airport in Lagos aboard KLM Flight 588, a Boeing
777, at 11 on Christmas Eve and arrived at Schiphol Airport in Amsterdam a little early, at 5:37
a.m. on Christmas Day.

Three hours later, at 8:54 a.m., Northwest 253, an Airbus A330, took off for Detroit Metropolitan
Wayne County Airport, with three pilots, eight flight attendants and the 278 passengers.

Amsterdam has long been an airport of concern for American aviation security officials, like other
major gateways in Europe, including London, Brussels and Frankfurt, where the Transportation
Security Administration sees an unusually large number of hits from people on so-called selectee or
no-fly lists associated with security threats, one former senior Homeland Security official said.

In 2007, the Amsterdam airport began testing body-scanning machines that can find threats hidden
under passengers' clothing, but there are only 10 such machines out of 200 security checkpoints at
the sprawling airport. In the United States, the T.S.A. has begun to substitute similar machines,
called millimeter-wave technology, for walk-through metal detectors.

``Those will pick up anything underneath clothing,'' said Edmund S.~Hawley, who served as the
agency's administrator until January. ``If he had it taped to his leg, it could have easily
identified something there.''

Mr.~Hawley said of Al Qaeda and like-minded militants: ``They have been trying since 2001, and they
are going to keep trying. You have to keep your vigilance up over the long term. That is the hard
thing.''

\section{For Airline Passengers, Pat-Downs, Searches and Restroom Monitors}

\lettrine{A}{s}\mycalendar{Dec.'09}{28} Detroit's airport was rattled on Sunday by a second
frightening incident in three days, passengers at airports in the United States and around the world
encountered stiff layers of extra security, with international travelers undergoing newly required
bag inspections, body searches and questioning at security checkpoints and before they boarded
planes.

At the Detroit Metropolitan Wayne County Airport here, officials detained a passenger who caused a
disruption aboard Northwest Airlines Flight 235 from Amsterdam, the same flight involved in a
terrorism attempt on Friday, when a Nigerian man caused a fire by injecting chemicals into a device
taped to his leg.

Pilots on Sunday declared an emergency after a second man, also a Nigerian, spent an unusually long
time in an airplane restroom, said a Department of Homeland Security spokeswoman, Sara Kuban.

In a similar scene to the one Friday, the plane was taken to a remote corner of the airport, far
from the terminals, where it was surrounded by emergency vehicles. All of the baggage was removed
from the plane, lined up on the tarmac and searched by explosive-sniffing dogs.

The passenger, whose name was not made public, ``was removed from the flight and interviewed'' by
the Federal Bureau of Investigation.

``Indications at this time are that the individual's behavior is due to legitimate illness,''
Ms.~Kuban said, ``and no other suspicious behavior or materials have been found.''.

The incident came amid security measures that were begun by the Transportation Security
Administration in response to Friday's terrorism attempt.

The T.S.A. issued an update on its Web site Sunday that said passengers would be subject to greater
security, but its information was not as detailed as the memorandums sent by the agency to airlines
this weekend. The airlines said the new T.S.A. measures required an additional round of searches,
including body pat-downs at airport gates overseas.

International travelers were also told that they could not leave their seats for the last hour of a
flight, during which time they also could not use a pillow or blanket. They were also limited to one
piece of carry-on baggage, including a purse or briefcase, and that piece had to be stowed in an
overhead compartment for the last hour of a flight.

Airlines were ordered to turn off in-flight entertainment systems with maps showing a plane's
location, and pilots and flight crews were told not to make comments about cities or landmarks below
the flight path.

There also were unspecified measures at airports in the United States, where lines at screening
machines grew long. At the Detroit airport, officers in bright blue vests marked ``Police'' walked
through the check-in lobby.

Jodi Syens, of Holland, Mich., her husband, Marvin, and daughter Rachel were aboard a flight from
London that arrived in Detroit several hours behind schedule Sunday. The Syens, who had arrived at
Heathrow Airport three-and-a-half hours before their flight, said passengers were taken by their
aircraft rows up to the gate, where carry-on bags were thoroughly checked and the travelers were
patted down.

``It was long, but we were appreciative,'' Mrs.~Syens said.

Tightened security at Narita International Airport in Tokyo came as a surprise to Wen-Lung Huang, of
Ann Arbor, Mich., who traveled from Taipei with his wife, Linda, and infant daughter, Catherine. The
Huangs, who had two carts loaded with baggage, had to check Catherine's car seat and portable
stroller. Japanese airport officials wrapped each in plastic and heavy tape.

Mr.~Huang was skeptical of the new measures, saying, ``I don't know that storing everything and not
going to the lavatory for the last hour is going to help.''

The one-bag limit proved a headache for Nicole Williams, a Canadian who lives in Manhattan and was
headed home Sunday from Ottawa. She was busily stuffing a BlackBerry and other valuables from her
purse into a carry-on bag that also held her Boston terrier.

Ms.~Williams, who was flying to La Guardia Airport, said that because Air Canada had charged her
extra to take her dog in the cabin, she had mistakenly thought that the dog's bag was not covered by
the new carry-on restrictions.

There was a good deal of confusion among international passengers interviewed on arrival at several
American airports Sunday, with many saying they had been unaware of Friday's terrorism attempt and
had not been told of the new security measures until they were actually confronted with them.

Henry Chen, 48, a businessman who lives in San Francisco, said he was shocked to have a female
flight attendant barge in on him in the restroom while he was washing his face during a flight from
Seoul. ``It was kind of weird, to have a lady try to get in,'' he said. ``She said that they had to
watch people being in the restroom too long.''

Joel Barnes and Bryan Duncan, both 27, were sitting in a Starbucks at Los Angeles International
Airport, after a 13-hour flight from Brisbane, Australia. They were awaiting the arrival of a friend
who had sent them a text message to alert them that his flight from Vancouver had been delayed for
two hours because of heightened security measures at the airport there.

The men said their armpits and shoes had been searched before they boarded their flight, and
Mr.~Barnes said, ``They rubbed their hands on the soles of our feet.''

They also recounted how an hour before landing an announcement had been made that no one could get
up for the remainder of the flight.

``It was kind of funny,'' Mr.~Barnes said, ``because the previous announcement had been about the
danger of deep-vein thrombosis or strombosis or whatever you get from sitting for too long. We
laughed.''

Priya Prasad, 32, an administrative assistant who lives in Oakland, Calif., said she was annoyed by
the extra hour it took her to get through security when she boarded a flight in Mumbai. ``They're
being extra cautious, which I guess is fine,'' she said. ``But I don't understand what it is they're
looking for. They went through my bag three times, and still I got my scissors and tweezers on the
plane.''

Lee Hung Kyu, 40, who has worked as a Korean Air flight attendant for 15 years, said the new rules
posed challenges for flight crews. On a flight to San Francisco, a woman who claimed to be sick
protested when asked to give up her blanket an hour before landing.

Mr.~Lee, who had made two in-flight announcements warning of the blanket collection, said he could
not turn up the temperature on the plane, for fear of overheating the other passengers. ``I allowed
her to keep her blanket with one condition,'' he said. ``I checked it again to make sure there was
nothing between the blanket and her lap.''

For the most part, crew members said passengers were cooperative. But Mee Hyum Koo, a passenger on a
Korean Air flight from Seoul to New York, said the extra steps added to the anxiety of flying. ``I
don't feel good,'' she said. ``It's uncomfortable. Scary.''

Still, Sarah Woodhouse, of Norwich, England, said she could understand the new steps. Before leaving
Amsterdam for Newark Liberty International Airport, she said she was asked to switch on her camera
to prove it was not an explosive device.

``Everyone just accepted that that's what you have to do,'' Ms.~Woodhouse said. Paul Bidwell, her
traveling companion and a fellow high school teacher, said: ``I'm quite happy for them to do it.
It's peace of mind for everyone.''

There was little peace of mind among passengers landing here Sunday on Northwest Flight 253 from
Amsterdam, many of whom said they feared they were about to be victims of an attack similar to the
attempt made on the same flight as it landed Friday.

Ali Hashemian, a passenger from Fort Lauderdale, Fla., said a man who had been sitting in the last
row of the plane had been in the restroom for a long time when an air marshal, his badge visible,
began knocking on the door, but received no answer.

Mr.~Hashemian and other passengers said it was unclear to them what happened next, but when the
plane landed, the authorities escorted the man off. ``He was very quiet,'' Mr.~Hashemian said. ``He
knew he was being arrested.'' The Associated Press said he had been released after questioning.

Another passenger, Denise Kabalka-Chesney, of Romeo, Mich., said passengers were told that the man
had been ``unruly.''

The plane made an abrupt landing and then taxied to the remote location.

``When we made the touchdown, we knew something was going on,'' said another passenger, Hitesh
Desai. ``The captain asked us to be quiet and stay in our seats. It all felt a little surreal.''

Passengers eventually boarded buses for the international terminal, where they waited another 30
minutes for the baggage before departing around 5:15 p.m., more than four hours after the plane had
been scheduled to land.

``In the end,'' said Joydeep Yadav, of Rochester Hills, Mich., ``it was just an inconvenience.''

\section{In Las Vegas, Sports Books in a Pocket}

\lettrine{F}{or}\mycalendar{Dec.'09}{28} decades, sports bettors in this city have queued up at
sports book windows and handed over cash in exchange for paper tickets formalizing their bets.

But on a recent Sunday in the sports book at the M Resort in nearby Henderson, gamblers seeking
action on professional and college football games were engaging in a much different ritual: betting
through hand-held devices no larger than a smartphone.

The technology, part of a new system from Cantor Gaming, a subsidiary of the financial services firm
Cantor Fitzgerald, lets gamblers place bets on certain sporting events like football or baseball,
tennis matches and horse races from anywhere in the casino and play other games like blackjack and
baccarat.

Gamblers can also for the first time wager on the outcomes of events as the events transpire. When a
football team lines up for a field goal, for example, bettors can bet on whether the kicker will
make it or miss.

Not surprisingly, casino operators like it. They are always trying to make their books more
profitable. While annual casino revenue in Nevada hovers around \$10 billion, last year revenue from
sports bets was only \$125 million, or 1.25 percent of that amount.

``All of a sudden, these same people who were betting once or twice a game at the beginning of a
sporting event can place wagers every minute if they want,'' says Anthony A.~Marnell III, chief
executive of the M Resort. ``Having this technology changes the entire equation for everyone
involved.''

To wager, bettors still fork over cash to human ticket writers in the sports book, but those workers
simply convert the money into electronic credits, which are stored in a triple-encrypted database
hosted on backroom servers.

At the M Resort, gamblers have two options to use these credits: they can sit in front of
touch-screen monitors in the sports book that resemble airline check-in kiosks, or they can check
out location-aware wireless hand-helds, called eDecks. These devices then connect to the casino
servers over a secured wireless Internet connection.

The eDecks resemble chubby iPhones, and can be taken almost anywhere in the casino -- into the
lounges and restaurants, for example. Bettors, however, cannot use them (or any other devices) while
playing table games.

At the sports books at the Venetian and the Palazzo on the Strip, eDecks are the only option, and
have been rebranded as PocketCasino.

In both introductions, the portable devices offer standard sports bets and are already loaded with
software that allows supplemental wagering in the form of blackjack and baccarat games. The games
include basic wagers, as well as propositional bets that dynamically calculate odds based on the
cards dealt.

The betting during events option is called in-running betting, and Lee Amaitis, chief executive of
Cantor Gaming, says oddsmakers have dreamed of this for decades but have lacked the number-crunching
resources to hang lines at game speed.

To solve this problem, Cantor Gaming sets its in-game lines based upon algorithms from the financial
services world. Mr.~Amaitis says these algorithms are variations of those created by a sister
company, Cantor Index, a financial spread-betting company in London that offers bets on various
markets including equities, indices, bonds and commodities.

``We've created an environment in which we're trading sports,'' Mr.~Amaitis said. ``With markets,
realities and probabilities of future events are changing constantly. It's no different in sports;
each game has millions of permutations and potential outcomes we can offer as odds.''

Despite its pedigree, Cantor Gaming's eDeck technology had to overcome a number of challenges before
it was approved by the Nevada Gaming Commission in 2008. First was the issue of security. The
company requires all eDeck users to obtain a plastic card embedded with a Radio Frequency
Identification chip. When bettors check out a device, ticket writers program it to work only within
a few feet of that chip, so no other people can pick up the device and wager on the original user's
account.

Sunny Tara, Cantor Gaming's chief technology officer, says that as soon as the device can no longer
read the RFID chip, it prohibits any further betting before being reset.

Mr.~Tara said another security precaution was to protect all monetary transactions by relegating
them to back-end servers. ``The portable devices are sophisticated, but they really are just tools
to access the real technology, which sits on the servers,'' Mr.~Tara said.

He added that the devices transmitted data over channels ``not normally used'' for mainstream Wi-Fi.

A second challenge for Cantor Gaming has been persuading old-school sports bettors to embrace the
new approach.

Mark Goldman, director of race and sports for the Venetian and Palazzo Las Vegas, said his sports
books were licensing 200 eDecks. Older gamblers have been slow to adopt them, he said, but younger
bettors have shown interest. The overall betting has increased since October.

``Once most people know this is available, they seem pretty interested to try it out,'' Mr.~Goldman
said. ``This is something we'll probably have in our sports books for a long, long time.''

\section{Adding Fees and Fences on Media Sites}

\lettrine{O}{ver}\mycalendar{Dec.'09}{28} more than a decade, consumers became accustomed to the
sweet, steady flow of free news, pictures, videos and music on the Internet. Paying was for suckers
and old fogeys. Content, like wild horses, wanted to be free.

Now, however, there are growing signs that this free ride is drawing to a close.

Newspapers, including this one, are weighing whether to ask online readers to pay for at least some
of what they offer, as a handful of papers, like The Wall Street Journal and The Financial Times,
already do. Indeed, in the next several weeks, industry executives and analysts expect some
publications to take the plunge.

Rupert Murdoch, beyond charging for access to The Journal, has talked about forming a partnership
with a single search engine, which would pay him for the rights to scour the news and entertainment
programming produced by his company, the News Corporation, rather than letting all search engines
crawl his sites. Also Hulu, which is owned partly by Mr.~Murdoch's company, is considering charging
viewers to watch some of the TV shows it now streams free.

Magazine publishers, meanwhile, have banded together to try to create their own version of the
iTunes store, aiming for a day when they can sell enhanced versions of what they have been giving
away. And more and more media companies are planning to charge for apps on iPhones and other mobile
devices, as well as on the Amazon Kindle and other e-readers.

Media companies of all stripes built their business models on the assumption that advertising would
continue to pour into their coffers. But with advertising in a tailspin, they now must shrink, shut
down or find some way to shift more of the cost burden to consumers -- the same consumers who have
so blissfully become accustomed to Web content that costs nothing.

So will future consumers look back on 2010 as the year they finally had to reach into their own
pockets?

Industry experts have their doubts, saying that pay systems might work, but in limited ways and only
for some sites. Publishers who sounded early this year as though they were raring to go have not yet
taken the leap, and the executives who advocate change tend to range from vague to cautious in
making any predictions about fundamentally changing the finances of their battered businesses.

But one thing clearly has shifted already, in a year rife with magazine closures and newspaper
bankruptcies: conventional wisdom among media companies has swung hard from the belief that pay
walls would only curb traffic and stifle ad revenue, to the view that media businesses need to try
something new, because the current path appears to lead to extinction.

``Content providers see that the idea that everything has to be free, supported by ads, isn't
working well, and they're trying to put the toothpaste back into the tube, but only partially,''
said Alan D.~Mutter, a media consultant and blogger who has been an executive at digital media
companies.

He went on: ``So we're looking at some sort of an inflection point, at least in attitude. But I
haven't seen much realistic, hard-headed thinking about how that's going to happen, so I don't know
how much is really going to change.''

Ann S.~Moore, the chief executive of Time Inc., the nation's largest magazine publisher, said, ``A
lot is going to change over the next two years.'' But she conceded that it was very hard to predict
the shape of that change, and she said that adding pay walls alone probably would not work.

Of course, it is the established media, with their legacy of high operating costs and outdated
technology, that face this problem. Leaner, newer online competitors will continue to be free,
avidly picking up the users lost by sites that begin to charge.

Arianna Huffington, co-founder and editor in chief of The Huffington Post, predicted that much of
the talk of media's mining the Web for new revenue would never become reality -- and that if it did,
free sites like hers would benefit. Some of the plans now being laid might work, she said, but many
of them would just alienate the Internet users who click from one site to another, wherever links
and their curiosity take them.

``I'm not minimizing the fact that there's a need to experiment with multiple new business models,''
she said. ``I just don't believe in ignoring the current realities.''

For more than a decade, media companies have hoped for a day when they could either control access
to their products online or at least put a price on them that a mass market would bear. But that day
has never come. What has changed is the level of threat they face, given the worst advertising
downturn in memory.

Since the infancy of the Web, there have been predictions that by making information more plentiful
and accessible, prices would be steadily driven down, with no bottom in sight. At first, it did not
seem to matter: Internet advertising grew at a breakneck pace, and traditional media thrived even as
the assumption of free content took root online.

But eventually, the rise of the Internet punished most media, starting with the music industry, in
the form of file-sharing. That history offers an object lesson. Despite the success of iTunes and
other pay services, illegal downloads remain common.

Print publications are suffering most now, but digital distribution has grown in importance for
broadcast television. Nearly all of its content is now available free online, as broadcast media
lose audience and advertising. Book publishers are also fighting the tide; Simon \& Schuster said
recently that it would delay the release of e-book versions of 35 big titles, like Karl Rove's
memoir and a Don DeLillo novel, fearing that the \$9.99 digital versions would eat into sales of
hardcover copies.

Cable television has been an exception, thriving on subscriber fees, but even there, executives fret
that consumers are disentangling themselves from their cable boxes, free to pick and choose
individual programs online and watch on their TVs. Jeffrey L.~Bewkes, the chairman and chief
executive of Time Warner, has advanced a plan that he calls TV Everywhere, which would allow paying
cable television subscribers to view shows online for no extra charge.

Similarly, Comcast started a service this month that gives subscribers to its broadband Internet and
digital cable services access to its cable programming on the Web.

These efforts are not about wringing extra dollars from the Web but about preserving the current
economics of the business.

``We're saying, since those payments you have made have found their way to the networks and through
distributors that give you the connection, that we want to have you be able to watch all those
networks on broadband,'' Mr.~Bewkes said recently at an investor conference in New York.

A leading evangelist for the coming of a new era is Rupert Murdoch, who has said he envisions a
not-too-distant day when all of the News Corporation's news properties, including Fox News Channel,
The Times of London and The New York Post, charge online. He and his executives have repeatedly
criticized search engines and news aggregators, saying it was ``theft'' to profit from publishers'
work.

The News Corporation has been shopping around an online payment software system -- so far without
much success -- in hopes of playing pied piper to other publishers, and it is a charter member of
the group of magazine publishers that have banded together, in a consortium announced this month.
And there have been talks about the possibility of Microsoft paying for the exclusive rights to have
its Bing search engine direct users to News Corporation sites.

``Quality content is not free,'' Mr.~Murdoch wrote in The Wall Street Journal on Dec.~8, days after
delivering a similar message at a Federal Trade Commission workshop. ``In the future, good
journalism will depend on the ability of a news organization to attract customers by providing news
and information they are willing to pay for.''

People who have studied the problem argue that charging online would work only if consumers were
offered a much-improved product with the convenience of access anywhere, on any digital device --
the core idea behind the magazine consortium and its planned online store.

By that standard, much of the talk of wringing more money from Internet users rings hollow, said Jay
Rosen, a professor of journalism at New York University and a prominent blogger on media subjects.
``People who really think we have to charge or the industry is sunk would be more persuasive if they
said at the same time we have to add more value than we've been adding,'' he said.

And, most industry experts agree, entertainment will be easier to charge for than news. It may be
hard to prevent free distribution of an episode of ``The Office'' or ``NCIS,'' but the product is
unique, with no substitute being created by someone else.

A small number of publications already charge for Internet access, including The Wall Street
Journal, The Financial Times, Newsday, Consumer Reports and The Arkansas Democrat-Gazette. But they
tend to be either specialty products or near-monopolies in local markets, and they generally do not
charge enough to fundamentally alter their profit pictures.

But for most general-interest news, any paid site would be competing with alternative versions of
the same articles, delivered by multiple free news sources.

``One of the problems is newspapers fired so many journalists and turned them loose to start so many
blogs,'' Mr.~Mutter said. ``They should have executed them. They wouldn't have had competition. But
they foolishly let them out alive.''

\section{Why Studios Keep Cranking Out TV Remakes, Despite the Flops}

\lettrine{I}{n}\mycalendar{Dec.'09}{28} the fall of 2000 all the buzz in the television business was
about a new drama on CBS on Friday night.

Not ``CSI'' -- that had no buzz at all. All it had was ratings. No, the buzz that fall was about the
remake of ``The Fugitive,'' the classic innocent-man-on-the-run series from the 1960s.

The redo seemed a can't-miss idea. But it lasted 23 episodes.

And then there was ``Bionic Woman.'' The NBC remake in 2007 of that sci-fi chestnut from the 1970s
started out as the hottest new show of that fall season -- and was gone after all of nine episodes.

There might be a lesson in there somewhere, but you would not know it from looking at the
development lists at three networks this winter. Among the most prominent projects under
consideration as new series next fall are these familiar names: ``The Rockford Files'' on NBC;
``Charlie's Angels'' on ABC; and ``Hawaii Five-O'' on CBS.

All of the projects were announced with some fanfare by their networks, though the program creators
and top network entertainment executives were reluctant to discuss any specifics about the new
versions yet, saying they were still in the writing stages. But the network executives expressed
genuine excitement about the possibilities for the projects.

It is easy to understand why. ``It's a good idea to try,'' said Warren Littlefield, who was the top
programmer at NBC and is now an independent producer. ``Movies have proved you can do well with a
presold concept.''

That is another way of saying it is only natural to turn to familiar titles because they attract
attention. The question is whether the series that result will attract viewers.

The track record does more than suggest not: it screams not. In the history of network television,
no remake of a previous hit series has ever become a hit itself on network television.

Plenty have been tried. In recent years there have been efforts to revive both ``Beverly Hills
90210'' and ``Melrose Place'' on the CW network. No one would claim either approaches the success of
their predecessors, or even passable hit status.

NBC created a splash when it brought back ``Knight Rider'' -- featuring a talking car -- as a
two-hour movie in 2008, but a series version later that year hit the skids quickly.

Earlier in the decade the highly successful ``Law \& Order'' producer Dick Wolf tried a remake of
the hoary classic ``Dragnet'' and barely made it through one season. (Who remembers that Ed O'Neill
played Joe Friday and that one of the co-stars was Eva Longoria?)

How many people could pass a quiz on who replaced Rod Serling as host of UPN's 2002 version of ``The
Twilight Zone''? (Forest Whitaker, actually.)

Digging through the files of series past, one could perhaps make an argument for some shows spawned
from original hits. ``Star Trek,'' of course, gave birth to four separate series, but those were all
spinoffs. They were not remakes of the original with the same characters.

All of those characters appeared instead in a series of theatrical movies, where television remakes
have been far more successful: ``Mission: Impossible,'' ``Get Smart,'' ``The Addams Family'' and
``The Brady Bunch,'' to name a few.

``Battlestar Galactica,'' a flop on ABC in 1978 (21 episodes), returned to a more favorable
reception on the Syfy network in 2004, but that was on cable, not broadcast television.

``Dragnet'' itself looks like a potential outlier. It started as a radio series and made it to
television in 1951. That run ended in 1959. But ``Dragnet'' came back in 1967 and was again a
success. That would seem to qualify it as a remake hit.

But the second show was less remake than revival: It still starred Jack Webb playing Joe Friday with
the same ominous theme music and intro about names being changed to protect the innocent.

The familiar titles on the current development slates are true remakes: ``Charlie's Angels'' will
star three new actresses as young women who are knockouts in looks and martial arts; ``Hawaii
Five-O'' is still expected to be led by a no-nonsense type named McGarrett (and maybe even backed by
a young aide named Danno); and ``The Rockford Files'' will inevitably feature a detective named Jim
Rockford.

Perhaps more important for that show is that it will also feature a creator with one of the best
recent r\'esum\'es in television. ``Rockford'' is being written by David Shore, the chief creative
force behind the hit drama ``House.'' Mr.~Shore has said that ``Rockford'' was one of his favorite
shows growing up and that he hopes to find a way to replicate its mixture of comedy and action.

But replicating a star on the level of James Garner, who played Rockford, may prove more
challenging.

Mr.~Littlefield said having a talent like Mr.~Shore running the show would be a great advantage to a
new ``Rockford,'' but he added, ``I don't think there are many gumshoe detectives around anymore, so
the key will be how they reinvent the character.''

``Hawaii Five-O'' is being created by the team of Alex Kurtzman and Roberto Orci, who wrote the
scripts for the most recent ``Star Trek'' movie as well as the two ``Transformers'' films.

The issue of how much to remake and how much to reinvent has dogged previous efforts at bringing
back familiar shows and characters. Fans and those who merely have heard of the old hits have tended
to turn up for the initial episodes (and for two hours' worth of a movie rendering), but have not
stayed around once they got a whiff of what the new version was really like.

``The identity of a hit TV series is so intimately tied to the original stars, style and attitude
that made it a hit in the first place that any deviation from that creates a real sense of aesthetic
dissonance,'' said Robert J.~Thompson, a professor of television and popular culture at Syracuse
University.

``This may be one case where an established brand is more a liability than an asset. In television,
it's a much safer approach to rip off an old idea than to try to remake one. It's a perfectly
plausible plan to develop a new TV show about three beautiful women fighting crime in fabulous
clothes; maybe not such a good idea to call it 'Charlie's Angels.' ''

Mr.~Littlefield said that the woeful track record of previous remakes should not discourage network
programmers from continuing to buy projects based on old hits. ``But there has to be a series
there,'' he said. ``It can't be like a movie. You can't trick them.''

Mr.~Littlefield suggested a formula that could work: ``At the risk of being oversimplistic: it also
has to be good.''

\section{A High-Tech Movie Battle: Which 3-D Glasses Are Best?}

\lettrine{W}{hile}\mycalendar{Dec.'09}{28} the blue-skinned Na'vi are shooting arrows out of the
screen toward the audience in the 3-D movie ``Avatar,'' another battle is being fought in the
theater -- over the goofy-looking glasses that moviegoers must wear to see the three-dimensional
effects.

Four companies are fighting for bridge of the nose with three different technologies. Each of them
is more advanced than the paper glasses worn to view ``Bwana Devil,'' regarded as the first of the
commercial 3-D movies in the 1950s, but all work on the same general principle. Each eye sees a
slightly different frame of the movie, but the brain puts them together and perceives depth.

About four million glasses made by RealD, the market leader, were worn during Avatar's opening
weekend in the United States. RealD's glasses use polarized lenses and cost about 65 cents each.
MasterImage 3D, another vendor, uses a similar technology.

Dolby Laboratories, the company behind theater sound systems, makes glasses that filter out
different frequencies of red, green and blue. They cost about \$28 each. The glasses of the third
company, XpanD, use battery-powered LCD shutters that open and shut so each eye sees the appropriate
frame of the movie. Those cost as much as \$50 each.

Each company claims its glasses and projection-system technology is better. Because glasses using
one technology are useless in a theater using a different digital projection system, the companies
backing the three technologies are scrambling for the upper hand while the 3-D industry is still in
its infancy.

James Cameron, the director of ``Avatar,'' is more often than not the main marketing tool. He has
endorsed RealD, says the company, which has about 5,000 screens using its system. But he, his wife
and his production partner were photographed at the premiere in Japan wearing XpanD glasses, which
work on 2,000 screens worldwide. Dolby says its glasses work with 2,200 screens, but it has no
Cameron connection. The company helpfully points out instead how a malfunction in the RealD system
spoiled a press preview of ``Avatar.''

The battle over what glasses patrons wear is a big deal because exhibitors are convinced that 3-D,
while seeming like a gimmick now, will lure movie lovers away from their crisp high-definition
widescreen TVs at home and back to the theater. But Maria Costeira, the chief executive of XpanD,
believes the sky's the limit: ``Eventually, we'll see 3-D movies on airplanes as well.''

The fight over the glasses may well intensify because TV makers are now pushing 3-D TVs for the home
as a way to increase their sales of more expensive sets.

Despite the marketing effort, when it comes down to choosing a 3-D system, many exhibitors are
making a decision based on one factor: Do they want to be in the cleaning as well as the movie
business?

The expensive Dolby and XpanD glasses are going into a dishwasher after each use, not the trash.
Both companies recommend that theater owners clean them in an industrial-grade machine. (To prevent
pilfering, Dolby and XpanD glasses can also contain built-in antitheft tags that can be activated by
exit-door sensors.)

XpanD offers its theater partners disposable wipes that it can distribute to customers along with
their tickets to assure them the glasses are germ-free.

RealD, whose cheap throwaway glasses were being perceived as a liability, has addressed concerns of
hygiene. Theater owners are now encouraged to ship back the used glasses to the company, which will
clean, repair and repackage them for other theaters.

But in all the hubbub about each product's advantages and which system Mr.~Cameron really, truly
loves, the most important question remains unanswered: does one system create a better looking 3-D
picture than another?

``I don't think the consumer can tell the difference,'' said Joe Miraglia, the director of design,
construction, and facilities for ArcLight Cinemas, a chain of luxury theaters based in Hollywood.
The movie chain uses each system in one or more of its theaters, and finds the cost of operation to
be roughly the same for all.

While Mr.~Miraglia uses RealD in several theaters, he chose XpanD's LCD glasses for the large curved
screen in the company's flagship Cinerama Dome theater on Sunset Boulevard. This is similar to the
technology that will be used by Panasonic, Sony and others as they bring 3-D HDTV to market next
year. Recently, electronics makers set standards for creating 3-D Blu-ray discs and players.

But in order to make the wearing of 3-D glasses as routine as ordering popcorn, the makers need some
help in the design department. Many of the glasses resemble the ``fitover,'' or wraparound
sunglasses favored by senior citizens in the Sun Belt, a look that is not appealing to young
moviegoers.

RealD and XpanD hope that 3-D will soon become a fashion statement. In addition to its standard
movie glasses, RealD is introducing child-size versions, as well as high-style 3-D specs that people
can wear without embarrassment out in the three-dimensional world as sunglasses or prescription
lenses.

Ms.~Costeira of XpanD thinks personalized designs that can be used with 3-D HDTVs and video games
could turn into something big. ``Stylish, thin and light, 3-D glasses will become your new iPod,''
she said.

\section{For Marketers, Love Is in the Air}

\lettrine{N}{ew}\mycalendar{Dec.'09}{28} commercials for BlackBerry feature young go-getters, like a
woman joining an all-male break-dancing group, a pair of designers sewing fabric for a triumphant
fashion show, and a fledgling rock band performing for an ecstatic audience. The company's
smartphones make only the briefest of cameos.

The real star, it turns out, is love.

``All You Need Is Love'' by the Beatles is the only audio in the spots -- by the Toronto office of
Leo Burnett, part of the Publicis Groupe -- and the ads culminate with the slogan ``Love what you
do.''

It is just one of a spate of campaigns focused on love, a sentiment enthralling Madison Avenue in
spite of -- or perhaps as an antidote to -- a downturn and two wars. Love is selling cars (``Love.
It's what makes a Subaru, a Subaru), LensCrafters eyeglasses (``See what you love, love what you
see.'') and Payless shoes (``I $[$heart$]$ shoes.''), not to mention long-running campaigns for
McDonald's (``I'm lovin' it.'') and Olay (``Love the skin you're in.'').

How the word for the most profound of human emotions came to be so popular for peddling consumer
goods has less to do with linguistics than psychology.

``There are left-brain and right-brain approaches to advertising,'' said Linda Kaplan Thaler, chief
executive of the Kaplan Thaler Group, part of the Publicis Groupe. ``For a long time, there was a
left-brain approach'' that highlighted ``rational reasons and good selling points, but in the last
several years studies have shown that emotional attachments really are a crucial factor in
purchasing decisions.''

Jerry Della Femina, chief executive of Della Femina Rothschild Jeary \& Partners, is dubious about
some love-themed ads.

``I wouldn't have the nerve to walk into a client's and say, 'Here's the campaign: everybody loves
your product,' '' said Mr.~Della Femina, who has worked on Madison Avenue since the 1960s.

But he is not surprised that such pitches are approved. ``Is a client going to say, when that's
presented, 'No, nobody loves us' ''? asked Mr.~Della Femina. ``Of course not -- they're going to
say, 'Everybody loves us.' ''

•

In an ad for Subaru -- by Carmichael Lynch in Minneapolis, part of the Interpublic Group -- a man
drives his Forrester for two days to ``Subaru Heaven,'' a junkyard where Subarus are left for parts.

``You don't just let some wrecker haul off your 300,000-mile Forrester to who knows where,'' says
the driver. ``You give that car a chance to live on -- one part at a time.''

The commercial is based on a letter sent by a customer, said Kevin Mayer, director of marketing
communication at Subaru, which introduced the ``It's what makes a Subaru, a Subaru'' campaign in
2007.

``We think of our customers as experience seekers, and the response they fed back to us over and
over is the love they had for the brand and how the product enables their lifestyles,'' Mr.~Mayer
said, adding that many Subaru owners are skiers or kayakers traveling unpredictable roads, and they
choose four-wheel-drive vehicles.

At Honda, a campaign that started in September -- by RPA in Santa Monica, Calif. -- features
nonactors describing Honda owners, closing with the slogan, ``Everybody knows somebody who loves a
Honda.''

Tom Peyton, senior manager for advertising at the American Honda Motor Company, said that as
domestic automakers have struggled this year, Honda, which has remained profitable, took a
softer-sell approach. ``It's been such a difficult year that the standard car ad you've seen out
there is some sort of incentive offer, some kind of lower-funnel, 'come buy it now for this reason'
approach,'' he said. ``Not too many companies have had the luxury of having just an upper-funnel
message that espouses love.''

A Nissan commercial that made its debut in October -- by the Los Angeles office of TBWA/Chiat/Day,
part of the TBWA Worldwide unit of the Omnicom Group -- shows an Altima enduring a pothole and a
student driver, to the Nazareth song ``Love Hurts.'' It closes with Nissan's new slogan, ``Quality
you can love.''

Asked how three automakers all arrived at love slogans, Erich Marx, director of Nissan marketing,
responded, ``Maybe we all felt there had been an emotional tug missing from the messaging.''

•

``The morning is jealous of your relationship with running,'' begins a recent New Balance
commercial, as a woman on the edge of her bed ties her laces, sheets magically wrapping themselves
around her, trying to draw her back into bed. `` 'What do you see in running anyway?' But somewhere
in the back of your mind is the thought of the perfect mile, politely telling the bed -- to shut
up.'' As the woman bounds from her front door, the voice-over concludes that with New Balance, ``you
can hate the morning less, and love running more.''

Norma Delaney, senior manager of global advertising and brand strategy for New Balance, stressed
that its ``Love/hate'' campaign -- by BBDO New York, part of the BBDO Worldwide unit of the Omnicom
Group -- is not boasting about lovability.

``If I look at some other work out in the marketplace, it's about telling the consumer to love the
brand more,'' Ms.~Delaney said. ``But we feel like the relevancy for our brand comes from
demonstrating that we understand what the consumer is going through -- this love-hate relationship
with running.''

Olay, the Procter \& Gamble brand, has used the ``Love the skin you're in'' slogan for about a
decade.

``We view Olay as a partner alongside women, so the emotional connection is Olay validating to a
woman that we want to help her achieve her best skin, to get to a place where she loves her skin,''
said Michael Kuremsky, vice president and global Olay brand franchise leader.

Research In Motion, parent company of BlackBerry, declined an interview request but provided a
statement explaining that ``Love what you do'' was inspired by ``the strong connection that people
feel with the BlackBerry brand, which is very often described and crystallized by BlackBerry
customers with the word 'love.' ''

What may be tempering BlackBerry users' love, however, is three widespread Internet and e-mail
lapses in December.

Although love cannot be measured precisely, of course, BrandIndex tracks consumer perceptions
through polling, subtracting negative responses from positive ones to assign companies a score from
negative 100 to 100.

When it comes to respondents' willingness to recommend a product, BlackBerry's score has dropped to
16, from 20, since the beginning of 2009, while Apple, the maker of the iPhone, has remained steady
at 36. Motorola, with its recent introduction of the Droid smartphone, has grown to 30, from 26.

(An ad for the Droid in The New York Times last week featured testimonials from users, who used the
word love six times, including this, from Kat Hubbard Amundsen: ``I love my Droid more than my kids,
more than my husband.'')

As for buzz, which measures positive news that respondents have heard about companies, Apple again
bests BlackBerry, with a 36 score consistent throughout the year, while BlackBerry slipped to 27,
from 34, and Motorola rose to 19, from 8.

``Buzz scores are probably the most sensitive to major advertising campaigns,'' said Ted Marzilli, a
managing director at BrandIndex. ``When I look at BlackBerry it seems pretty flat, and if you were
out there spending that kind of money you'd want more buzz.''

\section{Cellphone Companies Rush to Sue One Another}

\lettrine{W}{hen}\mycalendar{Dec.'09}{28} Nokia, the world's largest maker of mobile phones, sued
Apple, Samsung, LG and eight other competitors within six weeks beginning in October, it said it was
conducting a routine defense of its intellectual property.

But for cellphone makers and suppliers accustomed to swapping valuable technologies, the suits filed
by Nokia were far from standard.

Like many cellphone makers, Nokia is fighting the economic downturn. The company, which is based in
Finland, has laid off thousands of employees this year to counter falling sales and profit and its
slipping share of the global market, which fell to 35 percent in the third quarter, from 41 percent
in the second.

Bill Tai, a partner at Charles River Ventures, a technology investment firm, said the new legal
aggressiveness was ``a natural evolution'' similar to what took place in the semiconductor and
desktop computer businesses during difficult or competitive times.

``As competition has intensified and margins have deteriorated, companies are trying to put up
barriers to protect their positions,'' said Mr.~Tai, whose firm has invested about \$2 billion in
technology start-ups since 1970.

The mobile industry, since its inception, has been a legal battleground. In 2008, Nokia settled a
suit with Qualcomm, a United States maker of mobile phone chips, and agreed to pay an estimated
\$400 million a year for Qualcomm patents over 15 years. According to recent financial reports,
nearly all cellphone makers are suing or being sued.

Research In Motion, the maker of the BlackBerry, is suing Motorola, accusing it of using 20 of its
patents; Motorola has countersued. Palm is being sued by Saxon Innovations, a patent licensing
company in Tyler, Tex., over a third-party application processor. In March, Apple won a ruling in
Japan against Shigeru Saito Architecture Institute, which had sued over touch-screen patents.

``Nokia has been very active of late, and in general, there has been an increase in litigation in
the industry,'' said Clive D.~Thorne, an intellectual property lawyer at Arnold \& Porter in London.
``There is an attempt to gain technical advantage. Mobile technology is fast-moving. The commercial
life of products is limited. There is urgency in ensuring that I.P. rights are protected.''

In October, Nokia filed suit against Apple in United States District Court in Delaware, accusing the
company of violating 10 of Nokia's third-generation wireless patents with the iPhone. Then, in the
United States and Britain, Nokia accused Samsung, LG, Philips, Sharp and others of operating a
cartel to keep LCD screen prices high.

Louise Pentland, Nokia's chief legal officer, said the company's legal strategy had not changed.
What has, she said, is the dynamics of the mobile phone business: increasing boundary testing by
industry newcomers and suppliers, in addition to convergence of mobile Internet, information
technology and other industries.

``The litigation is not the driver,'' Ms.~Pentland said in an interview. ``It is the business
environment that has changed.''

Another reason for the rise in litigiousness is the increase in the number of patents sought by
makers of mobile phones and equipment in their relentless quest for the cutting edge, said Ian Drew,
an executive vice president at ARM, a company in Britain that has licensed its Cortex cellphone chip
designs to most of the world's cellphone makers and parts suppliers.

ARM's clients, in turn, use its designs for their own products, which they inevitably seek to
protect through patents. This so-called patent differentiation, Mr.~Drew said, of everything from
chips to touch-screen input methods, has created a legal landscape prone to conflict.

``In the mobile phone industry, there are lots of semiconductor and handset players jockeying for
position at the moment,'' Mr.~Drew said. ``Because there are so many smaller players, it is a more
dynamic business and you are going to get the odd skirmish.''

As cellphones become more like tiny computers, intellectual property will take on increasing
importance to equipment makers in their quest for market advantage.

For example, the Rosum Corporation, a technology company that developed a method for pinpointing the
physical location of mobile devices in dense urban areas with the aid of television broadcast
signals, is vigorously protecting its investment.

Rosum, which is privately owned and based in Sunnyvale, Calif., has only 24 employees but has
already obtained or applied for 71 patents.

``We have taken a great deal of care to protect our technology,'' said Todd Young, a vice president
for business development at Rosum. ``We think it is critical to our business.''

Such investments can be risky. When a company makes a patent infringement claim, it is often
countersued by its opponent. And often, the costs can be so high as to make the legal victories
almost worthless.

Nokia filed its suit against Apple more than two years after the iPhone went on sale, which
suggested that the suit was a symbolic line in the sand, said Jan Ihrfelt, an analyst at Swedbank,
one of the largest Swedish banks.

``There is no expectation that Nokia will gain a financial windfall from these lawsuits,'' he said.
``I think it was more Nokia trying to send a message.''

\section{Awkward Timing for a Book by Woods}

\lettrine{M}{otorists}\mycalendar{Dec.'09}{28} could soon be driving under the influence of Tiger
Woods, opening another round of possibilities for late-night talk show hosts.

Mr.~Woods has been the subject of an almost endless stream of jokes since crashing his sport utility
vehicle in his driveway, leading him to acknowledge his marital infidelities and take an indefinite
leave from professional golf. Among other things, his skills behind the wheel have been compared to
his precision when driving a golf ball.

And now, in what is not a joke, but might sound like a setup for one, Hachette Audio in May will
release Mr.~Woods's 2001 best seller, ``How I Play Golf,'' as an audio book. According to the
magazine AudioFile, 53 percent of audio book customers listen to the books while driving.

In the audio book, Mr.~Woods shares the ``psychological practices he uses daily to keep his game in
top shape and help him transcend all the ups and downs of golf,'' according to a description in
Hachette Audio's spring-summer 2010 catalog, which was mailed recently to bookstores and
journalists.

``The catalog went to bed months before the scandal unfolded,'' Anthony Goff, publisher and director
of Hachette Audio and Digital Media, wrote in an e-mail message. ``We had no idea he'd be all over
the news for anything other than his golf game.''

Mr.~Woods's favorability rating has plummeted recently to 33 percent, from a prescandal 85 percent,
according to a USA Today/Gallup Poll; advertisers including Accenture, Gillette and Tag Heuer have
either dropped him or removed ads that feature him.

The catalog says the book will be promoted with ``Father's Day promotions'' like retail display
racks -- awkward, perhaps, since the fact that Mr.~Woods has two young children has intensified the
disenchantment of some fans.

``We feel any marital infidelities wouldn't impact his solid golf advice in the least,'' Mr.~Goff
wrote. ``We were selling this title into stores for Father's Day promotions not for personal
relationship advice, but for tips on the game, which are still completely valuable to anyone who
knows golf.''

Mr.~Woods recorded his portions of the book (some segments are read by the actor Walter Franks) at
the time of the hardcover release in 2001, but the audio book was issued only on cassette, and has
not been available from the publisher for years. It will be sold on CD and as a digital download.

The audio release was originally scheduled to coincide with a paperback edition of ``How I Play
Golf,'' which has now been postponed, according to Rick Wolff, vice president and executive editor
of Grand Central Publishing, an imprint of Hachette.

The decision was made because the hardcover ``still sells well,'' Mr.~Wolff wrote in an e-mail
message. ``This was postponed long before Tiger ran into his current troubles.''

\section{China to Stay the Course on Currency, Wen Says}

\lettrine{P}{rime}\mycalendar{Dec.'09}{28} Minister Wen Jiabao of China struck a defiant note Sunday
about the country's exchange rate policy, saying the government would not give in to foreign demands
that it let the renminbi rise in value.

Mr.~Wen said in an interview with Xinhua, the official Chinese news agency, that the currency was
facing growing pressure to appreciate, but he insisted that China was committed to keeping it
stable, having virtually pegged it to the dollar since the global financial crisis worsened in the
middle of last year.

``We will not yield to any pressure of any form forcing us to appreciate,'' he said. As I have told
my foreign friends, on one hand, you are asking for the yuan to appreciate, and on the other hand,
you are taking all kinds of protectionist measures.'' The renminbi is known informally as the yuan.

The true purpose of these calls is to contain China's development, he said.

The renminbi has fallen against the currencies of most of China's trading partners this year because
it has been fixed to the weakening dollar, while China's economy has bounced back strongly.
U.S.~senators have asked for an investigation into whether current renminbi policy represents a
subsidy that would justify tariffs on Chinese imports.

Mr.~Wen also repeated an oft-made declaration that the stable renminbi had contributed to the global
economic recovery.

He gave a cautious outlook for the domestic economy in 2010, saying that it was too early to pare
down the government's stimulus policies but that officials needed to be attentive to surging real
estate prices and incipient inflation.

Although China will continue to encourage citizens to buy homes for their own use, differentiated
interest rates will be used as a tool to fight property market speculation, Mr.~Wen said.

He was apparently referring to a proposal that China keep offering preferential mortgages -- at a
discount of as much as 30 percent from benchmark lending rates -- for people buying their first
homes but eliminate such mortgages for additional home purchases.

More broadly, Mr.~Wen warned of rising imbalances from too much bank lending while defending China's
use of a stimulus package worth 4 trillion renminbi, or \$586 billion, to limit the effects of the
global economic crisis.

``Parts of the economy are not balanced, not coordinated and not sustainable,'' Mr.~Wen said,
repeating previous statements. He added that it would be better if lending by Chinese banks were not
on such a large scale.

China's overall lending situation improved in the second half of the year, when banks drastically
slowed their pace of credit issuance after a record surge in the first half, Mr.~Wen said. Chinese
banks are on course to lend an unprecedented 9.5 trillion renminbi this year, double the total of
the previous year.

\section{Bloggers Crash Fashion's Front Row}

\lettrine{N}{ot}\mycalendar{Dec.'09}{28} everyone thought it was adorable in September when a
13-year-old wunderkind blogger named Tavi was given a front-row seat at the fashion shows of Marc
Jacobs, Rodarte and others.

Oh now, don't misunderstand. She was totally adorable. You could have gobbled her up, with her goofy
spark plug style -- a Peggy Guggenheim for the Tweeting tween set. Her feet, in designer stockings,
did not quite touch the ground. Within a matter of months, Tavi Gevinson, the author of a blog
called Style Rookie, was feted by designers, filming promotions for Target, flown to Tokyo for a
party with the label Comme des Garçons and writing a review of the collections for no less than
Harper's Bazaar. Kate and Laura Mulleavy, the designers of Rodarte, described her in the pages of
Teen Vogue as ``curious and discerning.''

Rather, it was what the arrival of Ms.~Gevinson, as a blogger, represented that ruffled feathers
among the fashion elite. Anne Slowey, who has spent decades climbing the editorial ladder to a
senior position at Elle, dismissed the teenager's column as ``a bit gimmicky'' in an interview with
New York magazine. And in an instant, the subtext in her complaint was read by dozens of
Ms.~Gevinson's fans as an example of the tension between old media and new, when one leapfrogs ahead
of the other.

As a relatively new phenomenon in the crowded arena of journalists whose specialty it is to report
the news of the catwalks, fashion bloggers have ascended from the nosebleed seats to the front row
with such alacrity that a long-held social code among editors, one that prizes position and
experience above outward displays of ambition or enjoyment, has practically been obliterated. After
all, what is one to think -- besides publicity stunt -- when Bryan Boy, a pseudonymous,
style-obsessed blogger from the Philippines, is seated at the D \& G show in Milan between the
august front-row fixtures of Vogue and Vanity Fair, a mere two positions to the right of Anna
Wintour?

``There has been a complete change this year,'' said Kelly Cutrone, who has been organizing fashion
shows since 1987. ``Do I think, as a publicist, that I now have to have my eye on some kid who's
writing a blog in Oklahoma as much as I do on an editor from Vogue? Absolutely. Because once they
write something on the Internet, it's never coming down. And it's the first thing a designer is
going to see.''

Perhaps it was to be expected that the communications revolution would affect the makeup of the
fashion news media in much the same way it has changed the broader news media landscape. At a time
when magazines like Vogue, W, Glamour and Bazaar have pared their staffs and undergone deep cutbacks
because of the impact of the recession on their advertising sales, blogs have made remarkable
strides in gaining both readership and higher profiles. At the shows this year, there were more
seats reserved for editors from Fashionista, Fashionologie, Fashiontoast, Fashionair and others, and
fewer for reporters from regional newspapers that can no longer afford the expense of covering the
runways independently.

But it is somewhat surprising that designers are adjusting to the new breed of online reporter more
readily than magazines, which have been slow to adapt to the demand for instant content about all
things fashion. Blogs are posting images and reviews of collections before the last model exits the
runway, while magazine editors are still jockeying to feature those clothes in issues that will be
published months later.

So it is not without reason that some editors feel threatened, or that seasoned critics worry that
they could be replaced by a teenager. The designers and publicists who once quivered before the
mighty pens are now courting writers from Web sites that offer a direct pipeline to potential
customers. Sure, magazines and newspapers have started their own blogs and tweets, but reading them,
you often sense a generational disconnect, something like the queasy feeling of getting a ``friend''
request from your mother on Facebook. (From Glamour.com: ``Dating Tips: Why It's Important to Get
That Number.'')

Sites that include readers in the conversation are thriving, in a sense democratizing the coverage
of style, much as designers and retailers -- with lower priced fast-fashion collections -- have
democratized fashion itself. Garance Dor\'e and Scott Schuman, two photographers who have become
stars online (and who are a couple off-line), have created popular blogs with the simple idea of
posting images of stylish people and opening them to public comment. Now designers are seeking their
advice on communications strategies and even design -- Ms.~Dor\'e and Mr.~Schuman have worked on
projects with Gap.

Other sites have gained credibility along with traffic. Fashionista.com had 103,512 unique visitors
in November, and Fashionologie.com had 27,125, according to the online tracking agency Compete.
Jezebel.com (a saucy blog that includes coverage of fashion) shot ahead of Style.com (the Cond\'e
Nast fashion site) for the first time this fall with more than a half-million visitors. These are
considered large audiences for dispatches on such trivial developments as models refusing to wear
Alexander McQueen's crazy shoes or that such-and-such designer is looking for an intern.

The personalities behind those sites, in turn, are becoming as famous as some magazine editors. Marc
Jacobs named one of his bag designs after Bryan Boy, while Sephora asked Lauren Luke, whose makeup
videos are an Internet sensation, to preside over beauty contests in its stores. Designers are
thinking differently in response to consumers who want instant gratification. Doo-Ri Chung, for
example, describing her new basics collection in The Financial Times this fall, said her customer
has a ``blog mentality, not a magazine mentality.''

``The old idea of reading a magazine and planning ahead, that's not something that younger customers
do,'' she said. ``It's a different world, and designers have to adapt.''

Still, the popularity and novelty of such sites have raised concerns that their writers might be
unduly influenced by designers or beauty companies. New guidelines from the Federal Trade
Commission, announced in October, require blogs to disclose in their online product reviews if they
receive free merchandise or payment for the items they write about. This bothered some bloggers, and
reasonably so, since magazine editors commonly receive stockpiles of the same expensive goodies to
review in their pages, and that practice is rarely disclosed even though magazines are beholden to
advertisers for their livelihood.

Those guidelines also seem excessive at a time when magazines and newspapers are changing their tone
to embrace the online culture. On several fashion sites last week, it was reported that Vogue is
planning to feature a group of bloggers in its March issue, including Tommy Ton of Jak \& Jil, Ms.
Dor\'e and, yes, even Bryan Boy.

\section{Dissident Chinese Writer Appeals Sentence}

\lettrine{L}{iu}\mycalendar{Jan.'10}{05} Xiaobo, the Beijing writer who was given an 11-year prison
sentence on subversion charges on Dec.~25 after urging Chinese leaders to embrace democratic
reforms, appealed his conviction to the Beijing Supreme People's Court on Monday, according to one
of his lawyers.

The lawyer, Shang Baojun, told The Associated Press that Mr.~Liu told him during a meeting that he
had filed the appeal. ``We are making an effort to defend his innocence, but we have no idea whether
it will be successful,'' Mr.~Shang said.

Mr.~Liu's sentence, judged by many analysts to be unusually harsh, has drawn criticism from human
rights groups, Western governments and writers worldwide. Most regard the chances that it will be
overturned or softened as slim.

Mr.~Liu, 53, was seized by security officials in December 2008 as he and other intellectuals
prepared to issue Charter 08, a lengthy manifesto that called on China's Communist Party to uphold
individual rights and relinquish its monopoly on power. After being held more than a year in secret
detention and later in jail, Mr.~Liu was found guilty by a Beijing court on Dec.~25 of ``inciting
subversion of state power.''

More than 300 intellectuals and activists signed the charter, which was drafted in part by Mr.~Liu.
More than 10,000 others since have added their signatures to the document, which appeared briefly on
the Internet before Chinese censors banned it.

Legal experts said the 11-year prison term, the lengthiest handed down for subversion in at least a
decade, was intended to send a signal to other dissidents that Chinese leaders would tolerate no
challenge to the Communist Party's 60-year control of state power.

China's state-controlled news media have made no mention of Mr.~Liu's arrest, trial or sentencing.
The government news agency Xinhua issued a statement after Mr.~Liu's sentencing declaring that the
government had followed Chinese legal procedures in the case and that it had protected Mr.~Liu's
rights.

\section{For Shanghai Fair, Famous Fund-Raiser Delivers}

\lettrine{I}{n}\mycalendar{Jan.'10}{05} the hectic last week before she became secretary of state,
Hillary Rodham Clinton squeezed in a Bon Jovi benefit concert in New York, part of a frantic effort
to pay off the debt from her presidential campaign. No sooner had she arrived at the State
Department than Mrs.~Clinton discovered she needed to start raising money all over again.

This time, the cash-starved beneficiary was not her own campaign but the United States, which needed
\$61 million to finance the construction of a national pavilion at a world's fair in Shanghai. Under
federal law, no public money could be used for the project. And Mrs.~Clinton, as a federal official,
could no longer solicit private financial donations herself.

So she turned to her well-established network of Clinton fund-raisers, and after negotiating with
the State Department's lawyers about what she could legally do herself to support the project, she
mounted an ambitious fund-raising campaign that has netted close to \$54 million in barely nine
months.

With multimillion-dollar pledges from PepsiCo, General Electric, Chevron and other American
corporations, the United States is on track to open a sleek, 60,000-square-foot pavilion at the
Shanghai Expo 2010, which runs from May through October.

The prospect of the nation's chief diplomat asking for money worried government lawyers, according
to officials. Referring to the first secretary of state, one lawyer asked, ``Would Thomas Jefferson
do this?'' They imposed strict limits on the kinds of calls or other contacts she could make,
allowing her to promote the pavilion but prohibiting any one-on-one appeals for cash.

Despite those restrictions, and a dismal economy, Mrs.~Clinton is closing in on her \$61 million
goal. She is clearly proud of the effort, which staved off what could have been a rupture in
American-Chinese relations. In a year in which she has mostly worked to prove herself a loyal member
of the Obama team, the campaign also showcases her enduring political drawing power.

``The idea, for many people, of raising more than \$50 million would seem really daunting,''
Mrs.~Clinton said in an interview. ``Maybe because I had participated in raising so much money in
the past, I wasn't daunted by it. I knew it was going to be hard under the circumstances.''

By all accounts, the effort to build a national pavilion was near death at the end of the Bush
administration. The near-collapse of the global economy, the proximity of the expo to the Beijing
Olympics in 2008 and the general ambivalence of the State Department had left U.S.A.~Pavilion, the
nonprofit group in charge of the project, with little support or money.

``There is a sense in the U.S.~that Americans got disenchanted'' with world's fairs, said Nick
Winslow, a former Warner Brothers executive who is the president of U.S.A.~Pavilion.

With deadlines passing, the Chinese advanced the Americans money to conduct technical work for the
pavilion. They raised the issue with former President Jimmy Carter when he visited China last
January.

Enter Mrs.~Clinton, who made her first trip as secretary of state to Beijing in February and was
eager to talk about trade, climate change and the North Korean nuclear threat. Instead, she got an
earful about how bad it would be if the United States did not have a presence at the Shanghai Expo.

For the Chinese, the expo is a bookend to the Olympics. Shanghai is spending \$45 billion to
transform the city, even more than Beijing spent preparing for the Games. Nearly 200 countries have
signed on to take part, leaving only the United States and minuscule Andorra as potential no-shows.

``I was dumbfounded that so little attention had been paid to it,'' Mrs.~Clinton said. ``Everyone
knows China is going to be an enormously powerful player in the 21st century. They have an expo,
which is a kind of rite of passage that countries like to do to show they have arrived. We're not
there? What does that say?''

She said she did not relish the prospect of more fund-raising -- ``When would it ever end?'' she
recalled asking herself -- but she promised Chinese officials that she would try to raise the money.

There was little support within the State Department. So Mrs.~Clinton turned to two major
fund-raisers with long ties to the Clinton family: Elizabeth F.~Bagley and Jose H.~Villarreal.

Mrs.~Bagley, who is married to Smith Bagley, an heir to the R.~J.~Reynolds fortune, was ambassador
to Portugal under President Bill Clinton. Mrs.~Clinton appointed her to be the department's special
representative for global partnerships, a job that involves rounding up private support for public
projects.

Mr.~Villarreal, a well-connected San Antonio lawyer, has raised money for Mrs.~Clinton as well as
for Mr.~Clinton, former Vice President Al Gore and Senator John Kerry. In July, Mrs.~Clinton named
him the commissioner general to the expo.

To kick off the effort, Mrs.~Clinton held a conference call with 10 prominent chief executives.
Chevron, PepsiCo and General Electric each pledged \$5 million. Indra K.~Nooyi, the chief executive
of PepsiCo, made calls to other chief executives. Mrs.~Bagley and Mr.~Villarreal also opened their
Rolodexes, calling companies with operations in China. Some obvious prospects, like banks, were off
limits because they were receiving federal bailout money.

``In the beginning, we had to use a patriotism argument,'' said Kris M.~Balderston, Mrs.~Bagley's
deputy. ``The second wave of argument was commercial diplomacy. All of a sudden the companies
understood it would be good for them.''

Although Mrs.~Bagley is a State Department employee, she said she was advised that she could solicit
contributions. She noted that every would-be donor also had to be vetted by lawyers.

Fred Wertheimer, an advocate for stricter regulations for campaign fund-raising, said he was
satisfied that the State Department had handled a difficult situation properly.

``It would have been far better if the U.S.~government was able to pay for the activity involved,
but that does not appear to have been the case,'' he said.

While Mrs.~Clinton was barred from soliciting individuals, she met with corporate sponsors in
Shanghai in November, when she visited the expo site.

Her experience in the political trenches made a difference, Mr.~Villarreal said. ``Any other
diplomat would not have had the broad base of contacts,'' he said.

Mrs.~Clinton said it was easier raising funds for this project than to pay off campaign debt. ``I'm
much better at raising money for other people and other causes than I am for myself anyway,'' she
said, adding, ``Even though I've obviously raised a lot of money.''

Mark Landler reported from Washington, and David Barboza from Shanghai.

\section{For Some in Japan, Home Is a Tiny Plastic Bunk}

\lettrine{F}{or}\mycalendar{Jan.'10}{05} Atsushi Nakanishi, jobless since Christmas, home is a
cubicle barely bigger than a coffin -- one of dozens of berths stacked two units high in one of
central Tokyo's decrepit ``capsule'' hotels.

``It's just a place to crawl into and sleep,'' he said, rolling his neck and stroking his black suit
-- one of just two he owns after discarding the rest of his wardrobe for lack of space. ``You get
used to it.''

When Capsule Hotel Shinjuku 510 opened nearly two decades ago, Japan was just beginning to pull back
from its bubble economy, and the hotel's tiny plastic cubicles offered a night's refuge to salarymen
who had missed the last train home.

Now, Hotel Shinjuku 510's capsules, no larger than 6 1/2 feet long by 5 feet wide, and not tall
enough to stand up in, have become an affordable option for some people with nowhere else to go as
Japan endures its worst recession since World War II.

Once-booming exporters laid off workers en masse in 2009 as the global economic crisis pushed down
demand. Many of the newly unemployed, forced from their company-sponsored housing or unable to make
rent, have become homeless.

The country's woes have led the government to open emergency shelters over the New Year holiday in a
nationwide drive to help the homeless. The Democratic Party, which swept to power in September,
wants to avoid the fate of the previous pro-business government, which was caught off-guard when
unemployed workers pitched tents near public offices last year to call attention to their plight.

``In this bitter-cold New Year's season, the government intends to do all it can to help those who
face hardship,'' Prime Minister Yukio Hatoyama said in a video posted Dec.~26 on YouTube. ``You are
not alone.''

On Friday, he visited a Tokyo shelter housing 700 homeless people, telling reporters that ``help
can't wait.''

Mr.~Nakanishi considers himself relatively lucky. After working odd jobs on an Isuzu assembly line,
at pachinko parlors and as a security guard, Mr.~Nakanishi, 40, moved into the capsule hotel in
Tokyo's Shinjuku district in April to save on rent while he worked night shifts at a delivery
company.

Mr.~Nakanishi, who studied economics at a regional university, dreams of becoming a lawyer and pores
over legal manuals during the day. But with no job since Christmas, he does not know how much longer
he can afford a capsule bed.

The rent is surprisingly high for such a small space: 59,000 yen a month, or about \$640, for an
upper bunk. But with no upfront deposit or extra utility charges, and basic amenities like fresh
linens and free use of a communal bath and sauna, the cost is far less than renting an apartment in
Tokyo, Mr.~Nakanishi says.

Still, it is a bleak world where deep sleep is rare. The capsules do not have doors, only screens
that pull down. Every bump of the shoulder on the plastic walls, every muffled cough, echoes loudly
through the rows.

Each capsule is furnished only with a light, a small TV with earphones, coat hooks, a thin blanket
and a hard pillow of rice husks.

Most possessions, from shirts to shaving cream, must be kept in lockers. There is a common room with
old couches, a dining area and rows of sinks. Cigarette smoke is everywhere, as are security
cameras. But the hotel staff does its best to put guests at ease: ``Welcome home,'' employees say at
the entrance.

``Our main clients used to be salarymen who were out drinking and missed the last train,'' said
Tetsuya Akasako, head manager at the hotel.

But about two years ago, the hotel started to notice that guests were staying weeks, then months, he
said. This year, it introduced a reduced rent for dwellers of a month or longer; now, about 100 of
the hotel's 300 capsules are rented out by the month.

After requests from its long-term dwellers, the hotel received special government permission to let
them register their capsules as their official abode; that made it easier to land job interviews.

At 2 a.m. on one recent December night, two young women watched the American television show ``24''
on a TV inside the sauna. One said she had traveled to Tokyo from her native Gunma, north of the
city, to look for work. She intended to be a hostess at one of the capital's cabaret clubs, where
women engage in conversation with men for a fee.

The woman, 20, said she was hoping to land a job with a club that would put her up in an apartment.
She declined to give her name because she did not want her family to know her whereabouts.

``It's tough to live like this, but it won't be for too long,'' she said. ``At least there are more
jobs here than in Gunma.''

The government says about 15,800 people live on the streets in Japan, but aid groups put the figure
much higher, with at least 10,000 in Tokyo alone. Those numbers do not count the city's ``hidden''
homeless, like those who live in capsule hotels. There is also a floating population that sleeps
overnight in the country's many 24-hour Internet cafes and saunas.

The jobless rate, at 5.2 percent, is at a record high, and the number of households on welfare has
risen sharply. The country's 15.7 percent poverty rate is one of the highest among industrialized
nations.

These statistics have helped shatter an image, held since the country's rise as an industrial power
in the 1970s, that Japan is a classless society.

``When the country enjoyed rapid economic growth, standards of living improved across the board and
class differences were obscured,'' said Prof. Hiroshi Ishida of the University of Tokyo. ``With a
stagnating economy, class is more visible again.''

The government has poured money into bolstering Japan's social welfare system, promising cash
payments to households with children and abolishing tuition fees at public high schools.

Still, Naoto Iwaya, 46, is on the verge of joining the hopeless. A former tuna fisherman, he has
been living at another capsule hotel in Tokyo since August. He most recently worked on a landfill at
the city's Haneda Airport, but that job ended last month.

``I have looked and looked, but there are no jobs. Now my savings are almost gone,'' Mr.~Iwaya said,
after checking into an emergency shelter in Tokyo. He will be allowed to stay until Monday.

After that, he said, ``I don't know where I can go.''

\section{Large Oil Spill Reported in China}

\lettrine{A}{}\mycalendar{Jan.'10}{05} large oil spill in northwest China has heavily polluted a
tributary of the Yellow River, and threatens to reach one of the country's longest and most
important sources of water.

China's state-run news media said late Saturday that a ``large amount'' of diesel oil had leaked out
of a pipeline last Thursday in Shanxi Province.

The government has not explained why the report of the spill was not released until late Saturday.
But Xinhua, the official state news agency, said the leak was caused by construction work and that a
crew of 700 people was struggling to contain the damage from what Shanxi officials said was about
150,000 liters, or about 40,000 gallons, of diesel oil.

The damaged pipeline belongs to the China National Petroleum Corporation, one of the country's
state-owned oil giants and the parent company of PetroChina. The company did not specify on Saturday
exactly how much oil was spilled but said that it had shut down the pipeline. The company also said
that ``much of the leaked oil and polluted silt has already been taken away.'' But government
officials in Shanxi province said on Saturday that oil has been detected far downstream from the
leak and warned local residents not to use water in the region.

The oil pipeline, which transports oil from northwest China to central parts of the country, was
damaged and released oil into the Chisui River and Wei River, a tributary of the Yellow River,
according to Xinhua.

The Yellow River, which stretches for about 5,500 kilometers, or about 3,400 miles, is a source of
water for approximately 140 million residents, and it also provides water to factories and farms
through northern China.

In November 2005, a huge amount of toxic benzene leaked out and damaged the Songhua River in north
China's Heilongjiang Province, cutting off water supplies for millions of residents.

Some local officials were disciplined for a delay in reporting the spill, which later created a
panic among residents.

\section{Hong Kong Protesters Seek Democracy}

\lettrine{T}{housands}\mycalendar{Jan.'10}{05} of Hong Kong residents marched to the Chinese
government's liaison office on Friday demanding that Beijing grant full democracy to the
semiautonomous financial hub.

The turnout for the protest -- the police said 9,000 people took part -- was a boost to Hong Kong's
political opposition, which is trying to reignite the democracy movement at a time when locals are
more preoccupied with economic issues. At their peak, pro-democracy protests here drew hundreds of
thousands.

Five pro-democracy legislators plan to resign this month, hoping to turn the resulting special
elections into a referendum on democracy.

The former British colony returned to Chinese rule in 1997 under a separate political system that
promises Western-style civil liberties. The Chinese government ruled in 2007 that the territory
cannot directly elect its leader until 2017 and its legislature until 2020.

The protesters Friday said Beijing's timetable for democracy was too slow.

\section{Telecom Company to Pay \$3 Million in China Bribe Case}

\lettrine{U}{Tstarcom}\mycalendar{Jan.'10}{05}, an American telecommunications company, has agreed
to pay \$3 million in fines after United States law enforcement officials accused it of a
long-running scheme to bribe officials of China and other countries with cash, travel junkets and
other gifts.

The Justice Department and the Securities and Exchange Commission, which both issued statements on
Thursday, said that from 2002 to 2007, UTStarcom paid about \$7 million for hundreds of trips by
employees of China's government-owned telecommunications companies to visit the United States for
``training'' sessions that were often sightseeing vacations to Hawaii, Las Vegas and other tourist
locations.

The company also offered jobs to the family members of clients, paid them to attend universities,
secured their travel visas and arranged for a ``consultant'' to bribe a Mongolian official to win a
contract, the American officials said.

The actions appeared to violate the Foreign Corrupt Practices Act, which forbids bribing foreign
government officials, the officials said.

Those officials did not say why UTStarcom and its executives were not prosecuted.

Often, American officials choose not to prosecute because it is difficult to prove a case beyond a
reasonable doubt. Instead, they often opt to pressure companies to reform their operations and pay
large fines.

In a statement released on Thursday, UTStarcom said it took responsibility for the actions listed in
the government complaint. A company spokesman could not be reached for comment on Thursday.

The settlement is a window into how global companies sometimes do business in China, where bribery
and corruption are widespread.

Every year, tens of thousands of Chinese officials are arrested and accused of corruption. Many
legal experts say that steering clear of bribery here has become increasingly difficult, largely
because government officials control access to the country's booming economy.

Two years ago, Lucent Technologies, the American telecommunications giant, settled a similar case
with the S.E.C. after it was accused of spending \$10 million from 2000 to 2003 to pay for about
1,000 employees of Chinese government-owned telecommunications companies to inspect its plants and
attend training sessions. Instead, the groups traveled to Las Vegas, Disneyland, the Grand Canyon
and other tourist sites.

And just this week, China Mobile, one of China's biggest telecommunications companies, removed its
vice chairman because of ``serious economic crimes.''

For UTStarcom, which is based in Alameda, Calif., China is a crucial market. One of the company's
founders is a Chinese-born American, and most of the company's operations and employees are based in
China.

UTStarcom, which sells networking and broadband equipment, has for the last decade sold large
quantities of goods to three of China's biggest government-owned telecommunications companies: China
Netcom, China Telecom and China Mobile.

S.E.C. officials said employees from many of the company's big Chinese clients accepted bribes.

Apparently eager to maintain or increase its sales, UTStarcom engaged in complex frauds to please
its customers here, American officials said.

Company executives arranged for Chinese officials and their relatives to travel to the United States
and falsely recorded some of those trips as ``training sessions,'' even though no training took
place, the officials said.

The S.E.C. said there were also ``lavish gifts and all-expense-paid executive training sessions'' in
the United States for customers.

``This was apparently a standard practice,'' said Steven D.~Buchholz, a lawyer with the S.E.C.'s San
Francisco office.

Sometimes, the government said, customers were given cash allowances, and on other occasions,
UTStarcom put a Chinese official's relative on the payroll, even though that person never worked at
the company. The company ``paid and provided benefits to at least three of these individuals for a
period of two years each as if they were real employees, even though they never worked,'' the S.E.C.
said in its complaint.

The settlement comes at a difficult time for UTStarcom. The company is suffering through a sharp
downturn in sales. It lost over \$185 million in the first three quarters of 2008 and recently sold
one of its plants in China.

\section{Ex-Senator Runs for Governor in Rhode Island as Independent}

\lettrine{L}{incoln}\mycalendar{Jan.'10}{05} D.~Chafee, the former Republican senator from Rhode
Island who lost his seat to a Democrat in 2006 despite his family's longtime presence in state
politics, announced Monday that he would run for governor there as an independent.

Mr.~Chafee, 56, left the Republican Party after his loss to Sheldon Whitehouse, a Democrat, saying
the party's agenda had grown too conservative. In Congress, he opposed President George W.~Bush's
tax cuts and was the only Senate Republican to vote against authorizing the war in Iraq.

Gov.~Donald L.~Carcieri, a Republican who took office in 2003, cannot seek re-election because of
term limits, and the only Republican to enter the race dropped out last month. Mr.~Carcieri, a
conservative with flagging approval ratings, is the only Republican in statewide office.

On the Democratic side, Patrick C.~Lynch, the state's attorney general, and Frank T.~Caprio, its
general treasurer, are competing for their party's nomination. Both are well known in the state, but
so is Mr.~Chafee, whose family has loomed large in Rhode Island politics for more than a century.
His father, John, was the state's governor for six years in the 1960s before serving in the Senate
for two decades; upon his death in 1999, his son was appointed his successor.

``I believe that running as an independent will free me from the constraints that party politics
impose on candidates,'' Mr.~Chafee said in his announcement speech in Warwick, R.I. ``This freedom
will allow me to bring in the best people from both major parties and people without political ties
to solve our problems.''

During the recession, Rhode Island has suffered from some of the nation's worst budget problems and
one of its highest unemployment rates. In his speech, Mr.~Chafee said the state should ``carefully
examine'' expanding its sales tax to cover some items -- like food, clothing and over-the-counter
drugs -- that are currently exempt.

Mr.~Chafee faces significant challenges in the governor's race: Rhode Island is overwhelmingly
Democratic, and he trails his Democratic opponents in fund-raising. But New England voters have
elected a number of independents in recent decades, including two governors -- Lowell P.~Weicker
Jr.~in Connecticut and Angus King in Maine.

In Massachusetts, Timothy Cahill, the state treasurer, is running for governor as an independent
after leaving the Democratic Party last year.

Darrell M.~West, vice president of governance studies at the Brookings Institution, said Mr.~Chafee
would be a strong candidate even without the backing of a major political party. His name
recognition, his personal wealth and the fact that about half of Rhode Island voters are not
affiliated with a party all work in Mr.~Chafee's favor, Mr.~West said.

``The Chafee family name is the gold standard in Rhode Island,'' he said. ``And while Rhode Island
does tend to elect Democrats, independents are the largest voting bloc in the state.''

\section{Third Uninvited Guest at State Dinner}

\lettrine{T}{he}\mycalendar{Jan.'10}{05} saga of President Obama's first state dinner continues.

The Secret Service said Monday that a third uninvited guest gained entry to the dinner at the White
House on Nov.~24. A review of video from the party, which was held to honor the prime minister of
India, showed that a man wearing a tuxedo entered with members of the Indian delegation.

The discovery came as the Secret Service was investigating the case of Tareq and Michaele Salahi,
the Virginia couple who made their way into the White House as part of a publicity-seeking episode
involving a potential role in a reality television show. The Salahis came face to face with
Mr.~Obama, but officials said the third man, whose identity was not disclosed by the Secret Service,
apparently did not get near the president.

``At present, there is nothing to indicate that this individual went through the receiving line or
had contact with the president or first lady,'' Edwin Donovan, a spokesman for the Secret Service,
said in a statement Monday.

The Justice Department is investigating whether the Salahis violated the law by entering the White
House without an invitation. The case of the third man, which was first reported on the Web site
Newsmax, is also being reviewed by federal prosecutors to see whether he broke any laws.

The man is a United States citizen, a government official said, who was attending a meeting with
Indian leaders before the state dinner. He and others traveled from a Washington hotel to the White
House on a bus or van provided by the State Department.

``There was a group that was under our responsibility that went from a local hotel to the White
House,'' Ian Kelly, a State Department spokesman, told reporters. ``And there was a person who was
not authorized to be in that group who inserted himself or herself into that group.''

The incident has been reviewed, Mr.~Kelly said, and changes were made to how the State Department
handles visiting foreign delegations.

The uninvited guests at the dinner to honor Prime Minister Manmohan Singh of India, and his wife,
Gursharan Kaur, became one of the biggest controversies during the first year of the Obama
administration. Three uniformed Secret Service officers were placed on leave, and the White House
social secretary, Desir\'ee Rogers, was sharply criticized by members of Congress for her role in
allowing the Salahis to get into the White House.

The director of the Secret Service, Mark Sullivan, has acknowledged that the security breach should
not have happened. But he has said that neither the president nor the first lady was in any danger
during the dinner.

The Salahis, as well as the third man, passed through metal detectors and other screening before
entering the White House. But they did not undergo a background check, required for all visitors to
the White House, because their names were not on a list.

Telephone calls to the Indian Embassy in Washington were not returned on Monday. Neither the White
House nor the Secret Service would disclose the name of the third person who made his way into the
dinner.

\section{In Yemen, U.S.~Faces Leader Who Puts Family First}

\lettrine{T}{he}\mycalendar{Jan.'10}{05} United States is quickly ramping up its aid to Yemen, which
Washington sees as a revived new front against Al Qaeda. But one of the most delicate tasks will be
managing the relationship with the president of Yemen, Ali Abdullah Saleh, who has filled his
government with numerous members of his family and who wants to ensure that his son Ahmed succeeds
him, Yemeni officials, analysts and Western diplomats say.

Mr.~Saleh, 67, is wily, witty and fit. But he has been spending less time in the past two years
managing the complicated tribal and regional demands of fragile Yemen than trying to consolidate the
power of his family, the analysts say. As Yemen's oil revenues erode and Mr.~Saleh has fewer
resources to spread around, the reach of the central government has been shrinking -- ``the
government is practically caged in the capital,'' Sana, one senior Western diplomat said.

Mr.~Saleh presents the Obama administration with a problem that is all too familiar in Afghanistan
and Pakistan. He is amenable to American support, but his ineffective and corrupt bureaucracy has
limited reach. And his willingness to battle Al Qaeda, which he does not view as his main enemy, is
questionable.

Much of Yemen is in turmoil. Government forces on Monday killed two militants suspected of being
with Al Qaeda. There is another round of rebellion in the north and a growing secessionist movement
in the south. In important provinces where key oil resources are and where Al Qaeda in the Arabian
Peninsula is strong, government troops and the police largely remain in their barracks or in the
central cities. Order outside the cities is kept by tribal chiefs, with their own complicated
loyalties.

``You can't see anyone in a government uniform in Abyan,'' said Murad Zafir, a Yemeni political
analyst, referring to a southern province. ``There are large areas of the country where there is no
electricity, no running water and no central authority.''

United States aid was paltry until last year, and it was only when American intelligence could show
Mr.~Saleh that his family was being singled out by Al Qaeda that he began to take the group's threat
seriously, diplomats said.

How effectively Yemen addresses the threat depends largely on Mr.~Saleh's family. Ahmed Saleh is
head of the Yemen Republican Guard and the country's special forces.

The president's nephews -- sons of his late brother -- include Amar, the deputy director for
national security; Yahye, head of the central security forces and the counterterrorism unit; and
Tarek, head of the Presidential Guard. The president's half brother is head of the air force.

The sense of Yemen as a family corporation that has also enriched itself is part of the problem,
Mr.~Zafir said; the president's mosque, al-Saleh Mosque, was completed less than two years ago and
is said to have cost at least \$120 million. ``President Saleh wants his son to succeed him,''
Mr.~Zafir said. To make that happen, he has sought to consolidate power in his family's hands, but
his influence over the tribal chiefs has receded, Mr.~Zafir said.

Najeeb Saeed Ghanem, a former minister of health, is a member of Parliament from the largest
opposition party, Islah, an Islamist party with close ties to tribal groups. ``It is the size of the
deterioration of the regime and its control over the country that we're afraid of,'' he said.

With oil revenues down, Mr.~Saleh has had to turn to outside allies to help finance the war in the
north. Saudi Arabia provided \$2 billion last year to make up for the budget shortfall -- an amount
that dwarfs the \$150 million in security assistance that the United States will ask Congress to
approve for the 2010 fiscal year.

``The Saudis understand,'' said Ahmed M.~al-Kibsi, a political scientist at Sana University, ``that
they are the real prize for Al Qaeda, and Yemen is the platform.''

But there are challenges to Mr.~Saleh's goals of empowering his son, if not direct challenges to
Mr.~Saleh. One of his main allies, even as a young lieutenant colonel in 1977 when he initially took
power in the north, was Ali Mohsen. Now he is the military commander in charge of the effort to
stamp out the Houthi rebellion to the north. The Houthis are Shiites, and Mr.~Mohsen is said to be a
Sunni religious conservative.

Mr.~Saleh and Mr.~Mohsen are not related and are not considered rivals for the presidency. But
Mr.~Mohsen has signaled that he does not favor a direct succession of Ahmed Saleh to the presidency,
diplomats and analysts said. Mr.~Mohsen believes, they said, that the younger Mr.~Saleh lacks the
personal strength and charisma of his father and cannot hold the country together.

The tension between the two old comrades is visible in the criticism of the way the war in the north
is being handled, with government officials sometimes complaining that Mr.~Mohsen set off renewed
fighting there by occupying or destroying the mosques and holy places of the Houthis and building
Sunni mosques and schools in the area. Mr.~Mohsen's supporters have countered that the war has not
been fully supported by the central government.

Mr.~Saleh and his son also face another internal challenge from the next generation of the powerful
Ahmar family, Yemeni bluebloods. Sheik Abdullah al-Ahmar was the chief of the powerful Hashed tribe,
founded the Islah party and was Parliament speaker until his death in December 2007. One of his
sons, Hamid al-Ahmar, a businessman in his 40s, now leads Islah.

Mr.~Saleh has tried to keep the family close, in particular by letting Hamid al-Ahmar invest in
major cash cows, like the main cellphone company, SabaFon, oil interests and the Bank of Saba. But
in August, Hamid al-Ahmar stunned Yemenis by appearing on Al Jazeera to describe Mr.~Saleh as having
overstayed his time and calling on him to leave office and not try to enthrone his son. ``If Saleh
wants the people of Yemen to be on his side against monarchy and defend national unity, he himself
must quit pursuing monarchy,'' he said.

How the United States manages Mr.~Saleh and his family ambitions will have much to do with success
or failure against Al Qaeda. ``Washington must work with and behind the regime, whatever its flaws,
while trying to push Saleh toward reconciliation with his opponents,'' a Western diplomat said. ``I
am afraid it will take more delicacy than the Pentagon can do.''

\section{After Balmy Hawaii, Chilly Washington}

\lettrine{A}{s}\mycalendar{Jan.'10}{05} Air Force One lifted off late Sunday night, President Obama
and his family left behind a balmy 77-degree Hawaiian evening. Bearing east, the plane headed toward
Washington, where the temperature was a brisk 23 degrees.

If a 54-degree climate swing were not reason enough to lament the end of vacation, Mr.~Obama
certainly had others. He was returning to a Washington in the throes of a political furor over the
security breakdown that led to the attempted bombing of an American passenger jet on Christmas Day.

And if that were not daunting enough, the president faced a pile of other new and old business in
his Oval Office in-box. He has a State of the Union address to draft, a budget proposal to finish,
financial regulations to lobby for and a jobs plan to mull. Picking up where he left off just before
the holiday, he must also help fashion a compromise between competing House and Senate bills if he
wants to sign a health care overhaul into law in the next few weeks.

No wonder the Obamas seemed in a mood to rethink their departure. As they toured the Honolulu Zoo on
Sunday afternoon a few hours before packing up and breaking camp, a reporter asked the first family
if they were ready to go back to the mainland.

``No,'' they all cried in unison, led by the president.

``Let's stay,'' Michelle Obama volunteered. ``We'll all stay. Are we all in? I'm trying to mount a
coup.''

After all, this was not quite the vacation Mrs.~Obama had imagined. In the days leading up to the
holidays, she told visitors to the White House that she expected no news to break during the
family's 10-day trip to the president's home state and added that she would not be happy if any did.

Enter a Nigerian named Umar Farouk Abdulmutallab, who tried to ignite explosives hidden in his
underwear on Northwest Flight 253 as it approached Detroit. As a result, the president ended up
huddling each day with advisers who traveled here with him and conferring by secure telephone with
others back in Washington. He made two statements to reporters and taped a radio and Internet
address about the incident.

Talk of sticking around until Monday evaporated and instead the president opted for an overnight
flight so that he could arrive back at the White House by lunchtime on Monday. ``He's headed back on
a redeye so that he can sleep on the plane and jump right back into the full plate he's got in front
of him this week,'' said Bill Burton, a White House spokesman who accompanied the president.

No need to fret too much for the president. An overnight flight is not such an ordeal for him, since
he has a pretty decent bedroom aboard Air Force One. But staff members on long trips like that are
left to fend for themselves in seats or sometimes on pads laid out on the floor. At least no one
will stop them from going to the bathroom in the final hour of flight or order them to remove
blankets from their laps.

As much as the terrorist attack shifted the tone of the vacation, Mr.~Obama still managed to get in
a wide array of activity since arriving on Christmas Eve. He played golf and basketball with friends
like Eric Whitaker and Marty Nesbitt from Chicago as well as his old Hawaii pals, Bobby Titcomb,
Mike Ramos and Greg Orme. He went snorkeling with the family in picturesque Hanauma Bay. He took his
girls, Malia, 11, and Sasha, 8, to visit the local Sea Life Park, to see ``Avatar'' at a local
shopping mall movie theater cleared out just for them and to get the famous shaved ice at Island
Snow, where he ordered the ``Snowbama'' (lemon lime, cherry and passion guava).

For lunch, there were beach picnics plus an afternoon at Mr.~Titcomb's house on Oahu's North Shore.
Most evenings the family and friends, occupying three adjoining rental houses in Kailua, stayed in
for dinner, although they all headed out for a couple of excursions to Alan Wong's in Honolulu and
Lucy's Grill n' Bar near the compound in Kailua.

On his last day in Hawaii, Mr.~Obama paid a visit to the National Memorial Cemetery of the Pacific
at Punchbowl Crater, where his grandfather is entombed. There, as at most of his stops over the
course of the holiday, he made no public comments. In fact, for most of the trip, he remained out of
view of the pool of reporters and photographers that follows him everywhere he goes -- no pictures
of the presidential pecs this time.

But if Mr.~Obama still loves Hawaii, where he was born and spent most of his childhood and where he
has visited four times in the last 18 months, Hawaii certainly returns the favor. City and state
lawmakers are busy planning to rename all manner of places after the native-son president.

If his supporters have their way, according to the Honolulu Advertiser, there may soon be the
President Barack Obama Beach Park at Magic Island, the President Barack Obama Birthplace State Park,
President Barack Hussein Obama II Elementary School and President Barack Hussein Obama II High
School.

Moreover, the apartment building where he lived with his grandparents may be put on the National
Register of Historic Places. His birthday, Aug.~4, may be declared an official state holiday called
Barack Obama Day. And the anniversary of his inauguration, Jan.~20, may be recognized as Barack
Obama II Ohana Day. (Ohana refers to the sense of family in Hawaiian culture.)

Bathed in the love, it would not be all that surprising for Mr.~Obama to want to linger. But all
good things come to an end. After the longest break of his year-old presidency, Mr.~Obama had to
return to the cold, hard task of governing a nation still fighting two overseas wars and just
emerging from recession.

Sometime before Air Force One was to land at Andrews Air Force Base, Mr.~Obama presumably will have
shucked the khaki shorts and sunglasses in favor of a wool overcoat. And perhaps a coat of armor.

\section{How Visa, Using Card Fees, Dominates a Market}

\lettrine{E}{very}\mycalendar{Jan.'10}{05} day, millions of Americans stand at store checkout
counters and make a seemingly random decision: after swiping their debit card, they choose whether
to punch in a code, or to sign their name.

It is a pointless distinction to most consumers, since the price is the same either way. But behind
the scenes, billions of dollars are at stake.

When you sign a debit card receipt at a large retailer, the store pays your bank an average of 75
cents for every \$100 spent, more than twice as much as when you punch in a four-digit code.

The difference is so large that Costco will not allow you to sign for your debit purchase in its
checkout lines. Wal-Mart and Home Depot steer customers to use a PIN, the debit card norm outside
the United States.

Despite all this, signature debit cards dominate debit use in this country, accounting for 61
percent of all such transactions, even though PIN debit cards are less expensive and less vulnerable
to fraud.

How this came to be is largely a result of a successful if controversial strategy hatched decades
ago by Visa, the dominant payment network for credit and debit cards. It is an approach that has
benefited Visa and the nation's banks at the expense of merchants and, some argue, consumers.

Competition, of course, usually forces prices lower. But for payment networks like Visa and
MasterCard, competition in the card business is more about winning over banks that actually issue
the cards than consumers who use them. Visa and MasterCard set the fees that merchants must pay the
cardholder's bank. And higher fees mean higher profits for banks, even if it means that merchants
shift the cost to consumers.

Seizing on this odd twist, Visa enticed banks to embrace signature debit -- the higher-priced method
of handling debit cards -- and turned over the fees to banks as an incentive to issue more Visa
cards. At least initially, MasterCard and other rivals promoted PIN debit instead.

As debit cards became the preferred plastic in American wallets, Visa has turned its attention to
PIN debit too and increased its market share even more. And it has succeeded -- not by lowering the
fees that merchants pay, but often by pushing them up, making its bank customers happier.

In an effort to catch up, MasterCard and other rivals eventually raised fees on debit cards too,
sometimes higher than Visa, to try to woo bank customers back.

``What we witnessed was truly a perverse form of competition,'' said Ronald Congemi, the former
chief executive of Star Systems, one of the regional PIN-based networks that has struggled to
compete with Visa. ``They competed on the basis of raising prices. What other industry do you know
that gets away with that?''

Visa has managed to dominate the debit landscape despite more than a decade of litigation and
antitrust investigations into high fees and anticompetitive behavior, including a settlement in 2003
in which Visa paid \$2 billion that some predicted would inject more competition into the debit
industry.

Yet today, Visa has a commanding lead in signature debit in the United States, with a 73 percent
share. Its share of the domestic PIN debit market is smaller but growing, at 42 percent, making Visa
the biggest PIN network, according to The Nilson Report, an industry newsletter.

The Risk of Refusing

Critics complain that Visa does not fight fair, and that it used its market power to force merchants
to accept higher costs for debit cards. Merchants say they cannot refuse Visa cards because it would
result in lower sales.

``A dollar is no longer a dollar in this country,'' said Mallory Duncan, senior vice president of
the National Retail Federation, a trade association. ``It's a Visa dollar. It's only worth 99 cents
because they take a piece of every one.''

Visa officials say its critics are griping about debit products that have transformed the nation's
payment system, adding convenience for consumers and higher sales for merchants, while cutting the
hassle and expense of dealing with cash and checks. In recent years, New York cabbies and McDonald's
restaurants are among those reporting higher sales as a result of accepting plastic.

``At times we have a perspective problem,'' said William M.~Sheedy, Visa's president for the
Americas. ``Debit has become so mainstream, some of the people who have benefited have lost sight of
what their business model was, what their cost structure was.''

Visa officials said the costs of debit for merchants had not gone down because the cards now
provided greater value than they did five or 10 years ago. The costs must not be too onerous, they
say, because merchant acceptance has doubled in the last decade.

The fees are ``not a cost-based calculation, but a value-based calculation,'' said Elizabeth Buse,
Visa's global head of product.

As for Visa's market share, company officials maintain that it is rather small when considered
within the larger context of all payments, where, for now at least, cash remains king.

While Visa may be among the best-known brands in the world, how it operates is a mystery to many
consumers.

Visa does not distribute credit or debit cards, nor does it provide credit so consumers can buy
flat-screen televisions or a Starbucks latte. Those tasks are left to the banks, which owned Visa
until it went public in 2008.

Instead, Visa provides an electronic network that acts like a tollbooth, processing the transaction
between merchants and banks and collecting a fee that averages 5 or 6 cents every time. For the
financial year ended in June, Visa handled 40 billion transactions. Banks that issue Visa cards also
pay a separate licensing fee, based on payment volume. MasterCard, which is roughly half the size of
Visa, uses a similar model.

``It's a penny here or there,'' said Moshe Katri, an analyst who tracks the payments industry for
Cowen and Company. ``But when you have a billion transactions or more, it adds up.''

With debit transactions forecast to overtake cash purchases by 2012, the model has investors
swooning: Visa's stock traded at \$88.14 on Monday, near a 52-week high, while shares of MasterCard,
at \$256.84 each, have soared by more than 450 percent since the company went public in 2006.

While there is little controversy about the fees that Visa collects, some merchants are infuriated
by a separate, larger fee, called interchange, that Visa makes them pay each time a debit or credit
card is swiped. The fees, roughly 1 to 3 percent of each purchase, are forwarded to the cardholder's
bank to cover costs and promote the issuance of more Visa cards.

The banks have used interchange fees as a growing profit center and to pay for cardholder perks like
rewards programs. Interchange revenue has increased to \$45 billion today, from \$20 billion in
2002, driven in part by the surge in debit card use.

Some merchants say there should be no interchange fees on debit purchases, because the money comes
directly out of a checking account and does not include the risks and losses associated with credit
cards. Regardless, merchants say they inevitably pass on that cost to consumers; the National Retail
Federation says the interchange fees cost households an average of \$427 in 2008.

While the cost per transaction may seem small, at Best Buy, the biggest stand-alone electronics
chain, ``these skyrocketing fees add up to hundreds of millions of dollars every year,'' said Dee
O'Malley, director of financial services. ``Every additional dollar we are forced to pay credit card
companies is another dollar we can't use to hire employees, or pass along to our customers in the
form of savings.''

Weighing Rules on Merchants

The Justice Department is investigating if rules imposed by payment networks, including Visa, on
merchants regarding ``various payment forms'' are anticompetitive, a spokeswoman said. Several bills
have been introduced in Congress seeking to give merchants more ability to negotiate interchange,
which is largely unregulated.

While interchange remains legal despite repeated challenges, a group of merchants is pursuing yet
another class-action suit, this time in federal court in Brooklyn, against Visa and MasterCard that
seeks to upend the system for setting fees.

``Visa and MasterCard have morphed into a giant cookie jar for banks at the expense of consumers,''
said Mitch Goldstone, a plaintiff in the case.

Fees were not an issue when debit cards first gained traction in the 1980s. The small networks that
operated automated teller machines, like STAR, Pulse, MAC and NYCE, issued debit cards that required
a PIN. MasterCard had its own PIN debit network, called Maestro.

Merchants were not charged a fee for accepting PIN debit cards, and sometimes they even got a small
payment because it saved banks the cost of processing a paper check.

That changed after Visa entered the debit market. In the 1990s, Visa promoted a debit card that let
consumers access their checking account on the same network that processed its credit cards, which
required a signature.

To persuade the banks to issue more of its debit cards, Visa charged merchants for these
transactions and passed the money to the issuing banks. By 1999, Visa was setting fees of \$1.35 on
a \$100 purchase, while Maestro and other regional PIN networks charged less than a dime, Federal
Reserve data shows. Visa says the fee was justified because signature debit was so much more useful
than PIN debit; at the time, roughly 15 percent of merchants had keypads for entering a PIN.

Merchants said they had no choice but to continue taking the debit cards, despite the higher fees,
because Visa's rules required them to honor its debit cards if they chose to accept Visa's credit
cards.

A Seven-Year Battle

Wal-Mart, Circuit City, Sears and a number of major merchants eventually sued. After seven years of
litigation, Visa and MasterCard agreed to end the ``honor all cards'' rule between credit and debit
and to pay the retailers a settlement of around \$3 billion, one of the largest in American
corporate history. Visa paid \$2 billion, and MasterCard the remainder.

Since then, only a handful of retailers have stopped accepting Visa debit cards, an indication that
the crux of the lawsuit was ``much ado about nothing,'' Mr.~Sheedy says.

And while some merchants said they thought the lawsuit would pave the way to a new era of
competition, a curious thing happened instead: while Visa temporarily lowered its fees for signature
debit, it raised the price on PIN debit transactions and passed the funds on to card-issuing banks,
and its competitors soon followed.

The current class-action lawsuit joined by Mr.~Goldstone contends that Visa's PIN debit network,
called Interlink, is offering banks higher fees as an incentive to issue debit cards that are
exclusively routed over this network. Interlink, which has raised its PIN debit fees for small
merchants to 90 cents for each \$100 transaction, from 20 cents in 2002, is often the most
expensive, especially for small merchants, Fed data shows.

One large retailer, who requested anonymity to preserve its relationship with Visa, provided data
that showed Interlink's share of PIN purchases rose to 47 percent in 2009, from 20 percent in 2002,
even as its fees steadily increased ahead of most other networks -- to 49 cents per \$100
transaction in 2009, from 38 cents in 2006.

Visa officials say its PIN debit network is taking off despite rising costs because it offers
merchants, banks and consumers a level of efficiency and security that regional networks cannot
match. ``We are motivated as a company to try to drive value to each one of those participants so
that they accept the card, issue more cards, use the card,'' Mr.~Sheedy said.

At checkout counters, meanwhile, consumers are quietly tugged in one direction or the other.

Safeway, 7-Eleven and CVS drugstores automatically prompt consumers to do a less costly PIN debit
transaction. The banks, however, still steer consumers toward the more expensive form of signature
debit. Wells Fargo and Chase are among those that offer bonus points only on debit purchases
completed with a signature.

Visa says it does not care how consumers use their debit card, as long as it is a Visa. But for now
at least, the company says the only way to ensure that a purchase is routed over the Visa network is
to sign.

``When you use your Visa card, you have a chance to win a trip to the Olympic Winter Games,'' a new
Visa commercial promises.

The commercial does not explain the rules, but the fine print on Visa's Web site does: nearly all
Visa purchases are eligible -- as long as the cardholder does not enter a PIN.

\section{Divergent Views on Signs of Life in the Economy}

\lettrine{S}{igns}\mycalendar{Jan.'10}{05} of revival on the factory floor emerged in the United
States and Asia on Monday, fueling exultant buying on Wall Street and reinforcing a sense that
genuine economic recovery was unfolding.

But if the Great Recession has indeed relaxed its grip on American life, it has been replaced by
something that might be called the Great Ambiguity -- a time of considerable debate over the clarity
of economic indicators and the staying power of apparent improvements.

Manufacturing expanded in the United States in December, the fifth straight month of gains,
amplifying hopes that a job market hobbled by double-digit unemployment might finally be adding
paychecks. New jobless claims slipped markedly last week. Some economists think data to be released
on Friday will show the economy gained jobs in December, the first monthly net increase in two
years.

``We're really coming back,'' said Allen Sinai, chief global economist at the research firm Decision
Economics. ``The expansion is picking up the pace.''

Investors pushed up stock prices, with the Dow Jones industrial average gaining 156 points, or 1.5
percent, and oil prices reaching a 14-month high, \$81.51 a barrel. $[$Pages B1 and B8.$]$

But many economists remain worried that momentum could soon weaken, with the economy sliding back
into glum times.

Indeed, the only area in which economists can reliably declare expansion is in the supply of
competing narratives about the economy -- perhaps to be expected in any transition between downturn
and the inevitable turn for better.

``That is always the nature of the boomlet after recession,'' Mr.~Sinai said. ``People think it's
going to fade away.''

Much of the improvement in manufacturing -- a small slice of the American economy -- is the result
of businesses rebuilding inventories after slashing them. The economy has also been stoked by \$787
billion in federal spending aimed at stimulating growth, a force that will be largely exhausted by
the middle of the year.

Those skeptical of lasting recovery assert that, once businesses have rebuilt inventories and
federal largess runs dry, the economy will confront the same assortment of ills plaguing it for two
years.

Namely, Americans are saturated in debt and nervous about job prospects, prompting many to hunker
down in a mode of thrift; businesses still spooked by dysfunction in the financial system are
reluctant to hire more workers until recovery proves real; and a cataclysmic drop in home prices has
diminished spending power in millions of households, with another decline possible as foreclosed
properties surge onto the market.

``You come back to the unfortunate reality that the underlying problems are still very much there,''
said Joshua Shapiro, chief United States economist at the market research firm MFR Inc.~``It's going
to be a long grind.''

Few economists expect a double-dip recession, in which the economy begins to expand only to contract
again, but many fear a protracted period of stagnant growth that could last several years, much like
Japan's Lost Decade in the 1990s following the calamitous end to an era of excessive real estate
speculation.

Indeed, some liken signs of improvement in the American economy to events in Japan, where victory
was often prematurely declared.

``Japan had lots of upswings,'' said Stephen King, global chief economist at the banking
conglomerate HSBC in London.

Mr.~King foresees a long period of adjustment to slower rates of economic growth in the United
States.

The American economy expanded at a roughly 3 percent annual pace through most of the 1980s and
1990s, then dipped to about 2.5 percent from 2003 to 2007, and most of the gains stemmed from the
disastrous investment bubble in real estate.

``What are you left with in terms of underlying growth?'' he asked. ``The answer is, not very
much.''

But economists of a more optimistic type assert that such views fail to account for what has been
happening in recent months. Areas of the economy that went into deep freeze have revived and must
grow just to replenish their inventories.

Exhibit A in this narrative is the widely watched index of manufacturing activity produced by the
Institute for Supply Management, which reported on Monday increasingly robust orders and production
at American factories in December.

The purchasing manager's index rose to 55.9, its highest level since April 2006, back in the days
before subprime entered the American lexicon and before economists began to fret about a potential
replay of the Great Depression.Hopes for a lasting recovery were also bolstered by news from China,
where data released Monday showed its economy expanded in December at a faster clip than at any time
since the beginning of the global downturn. Similar data amplified evidence of recovery in South
Korea, India and Taiwan.

``Strong growth in Asia reverberates around the world and helps U.S.~exports,'' Mr.~Sinai said.

The worldwide factory revival is no momentary event, argues Robert J.~Barbera, chief economist at
the research and trading firm ITG, but rather the return to everyday business after an overreaction
to the terror that gripped financial markets in the fall of 2008.

During the panic, companies cut too many workers and pared back too much inventory in a frantic bid
to limit costs, Mr.~Barbera argues.

``You're having a global reversal of inventories and trade, and I don't think that's ephemeral,'' he
said. ``If it was important on the way down, why shouldn't it be important on the way back up?''

Mr.~Barbera anticipates the economy will add some 2.5 million new jobs this year -- far less than
the roughly 7 million net jobs that have been lost over the last two years, according to the Labor
Department, but still a considerable improvement.

As people spend paychecks at other businesses, they create jobs for more people, spreading the
improvement in national fortunes. This prospect captured attention on Wall Street on Monday. As
investors absorbed evidence of factory expansion, they saw signs of more cash registers ringing up
sales, more insurance agents writing new policies, more trucks carrying more product and more
transactions brokered by banks.

But whether growth in manufacturing indeed spells new jobs throughout the economy remains a
question. Mr.~Shapiro, the MFR economist, argues that companies have become permanently lean,
mastering the art of producing more with fewer people -- a trend with staying power.

A government report on Monday, for instance, showed a slump in construction spending, which fell a
seasonally adjusted 0.6 percent in November to its lowest levels in six years.

Stock markets may be reflecting the exuberance of improved news headlines more than fundamentals,
say analysts, and some buyers may be jumping in out of fear of missing the next fiesta.

``Typically, what happens with individual investors is they just chase the market,'' Mr.~Shapiro
said. ``There's an element of that for sure.''

But while pundits are prone to dismiss the stock market as a kind of indecisive judge at a beauty
pageant, its movements have substantial impact on wealth and confidence. These, in turn, shape
spending and employment.

In this context, the stock market may prove to be the wild card whose resilience cements a still
nascent recovery.

``It's part of a process by which households that were devastated are now improving their balance
sheets,'' Mr.~Sinai said. ``Their net worth is going up, and some of those households are going to
spend.''

\section{Google Moves to Keep Its Lead as Web Goes Mobile}

\lettrine{G}{oogle}\mycalendar{Jan.'10}{05}'s expected unveiling on Tuesday of a rival to the iPhone
is part of its careful plan to try to do what few other technology companies have done before:
retain its leadership as computing shifts from one generation to the next.

The rapid emergence of the smartphone as a versatile computing device may be as much a challenge as
an opportunity for Google, which built its multibillion-dollar empire largely on the sale of small
text ads linked to search queries typed on PCs.

As people increasingly rely on powerful mobile phones instead of PCs to access the Web, their
surfing habits are bound to change. What's more, online advertising could lose its role as the Web's
primary economic engine, putting Google's leadership role into question.

``The new paradigm is mobile computing and mobility,'' said David B.~Yoffie, a professor at the
Harvard Business School. ``That has the potential to change the economics of the Internet business
and to redistribute profits yet again.''

In recent decades, the power of industry giants like I.B.M. and Microsoft, which once seemed
unassailable, waned as computing shifted from big mainframes to PCs, and from PCs to the Internet.
Many analysts say it is now Google that is faced with a less certain future in the face of another
shift.

Still, they say Google saw this coming years ago and has been preparing for it. Google executives
now say they are confident that the company will thrive as the mobile Internet grows.

``We are incredibly excited about the opportunities that we see in mobile,'' Vic Gundotra, a vice
president of engineering at Google who oversees mobile applications, said in an interview on Monday.
``We have invested a considerable amount, and we can now really provide a compelling mobile
experience.''

Top Google executives, including Eric E.~Schmidt, the chief executive, have long said that the
mobile Internet was Google's biggest opportunity for new growth. They orchestrated a string of
acquisitions of companies with mobile-related technology, including Android, maker of a cellphone
operating system; GrandCentral, a service for making calls that can bypass telephone lines; and
AdMob, an advertising network for mobile applications. The AdMob deal is awaiting approval from
regulators.

Google also invested far more aggressively than its competitors in mapping technologies and services
tied to a user's location, which are likely to become the vital underpinnings of new advertising
systems on GPS-equipped mobile phones. Last month, Google came close to paying more than \$500
million to acquire Yelp, a Web site for business listings and reviews. While the deal collapsed at
the last minute, Google's interest underscored its determination to become a force in mobile
advertising.

And in recent years, Google has worked systematically to loosen the hold that other companies have
on the mobile industry.

In 2008, for example, Google bid \$4.7 billion in a government auction of the nation's airwaves.
While Google had no intention of winning, it bid to ensure that the airwaves would be subject to
so-called openness requirements, meaning that Verizon Wireless, which won the bidding, would not be
able to exclude Google services like Web search, Gmail and maps from phones using those frequencies.

The expected unveiling on Tuesday of the Nexus One, a thin, touch-screen handset built to Google's
specifications and made by the Taiwanese company HTC, is a challenge to a newly minted industry
power: Apple, whose iPhone dominates the high end of the smartphone market. While the iPhone sends
millions of people to Google's search and other services, some of the company's applications, like
Google Voice, have not been allowed to run on the phone.

Analysts say that with the Nexus One, which Google plans to sell to consumers directly, the company
is trying to free itself from Apple's growing influence. It also wants to broaden the appeal of
Android's technology. The phone is expected to be sold unlocked, allowing consumers to buy service
plans separately.

Mr.~Gundotra declined to discuss specifics of the Nexus One. But he said all of Google's mobile
moves were driven by one objective: pushing the industry to open up in an attempt to replicate on
mobile phones the environment that has allowed the PC-driven Web to grow at explosive rates.

``Before the mobile Web really started to take off, there were many barriers to consumers,'' he
said. ``Sometimes it was limited choice about what you could do with your phone,'' he said, adding
that in some cases, it could take as many as 19 clicks for a user to get to Google's site.

Some of Google's moves, like its bid for spectrum, confounded many in the industry. But analysts say
that Google's actions proved shrewd and that the company has, to a large extent, helped open up the
mobile Web and ensure that its own services, and its ads, will be accessible to all.

``You could take a view that this is a very geeky company,'' said Nicholas Carr, author of ``The Big
Switch,'' a book about the shift to Internet computing. ``That underestimates the strategy that
underlies all these moves.''

Some analysts say that with the early success of the Android operating system, which is built into
phones from several manufacturers, Google is already beating Microsoft, its biggest rival, in the
mobile business. And they note that mobile phone software is tethered to the Web more than PC
programs, playing to Google's strengths in Internet computing.

Indeed, Google has moved so aggressively to establish itself as a force in the mobile Web that it
has already attracted government scrutiny. The Federal Trade Commission recently stepped up its
review of Google's proposed \$750 million acquisition of AdMob, as some advocacy groups are raising
alarms that the deal could extend Google's dominance of online advertising into mobile phones.

Yet some questions remain unanswered: Will advertising remain central to the Web economy as
consumers shift to mobile phones from PCs? And will applications change people's reliance on search
engines?

Google's text ads are sold through an auction system, and analysts say that some of its most
lucrative ads show up during intensive online research tasks, like finding a vacation rental or
securing a good rate on a mortgage. Those are more likely to be conducted on a PC than on a
cellphone.

``It certainly remains to be seen how big mobile advertising will be,'' said Charles Golvin, an
analyst with Forrester Research.

Predictions about the growth potential of mobile advertising vary widely. A recent report on the
mobile Internet by Morgan Stanley, for instance, said that while advertising accounts for 40 percent
of revenue on the desktop Internet, it accounts for just 5 percent of revenue on the mobile
Internet. That could change, as more personalized advertising technologies, including coupons and
offers that are aimed at users based on their location, could usher in a new wave of growth in
digital marketing.

At the same time, some surveys show that users are wary of ads that could clutter the precious real
estate on their small cellphone screens. And phone users seem more willing to pay a few dollars for
applications or content than PC users, potentially reducing the importance of advertising.

Another risk for Google is that popular smartphone apps could erode the power of its search engine.
In a recent note to investors, Ben Schachter, an analyst with Broadpoint AmTech, wrote that apps
could eat into search revenue by giving users direct access to many commercial Web sites.

``When you go to Google and search for a product or for Amazon, that's a way for Google to make
money,'' Mr.~Schachter said. ``But if you have the Amazon or the eBay app on your mobile phone home
screen, you are more likely to click on that button and buy something without ever using Google.''

Mr.~Gundotra dismissed those concerns. While he refused to provide specifics about Google's mobile
revenue, he said that its rate of growth mirrored the company's results when it began selling ads on
its Web site.

``It is very, very encouraging,'' he said.

\section{A Venture Integrating Skype Into the Family Room}

\lettrine{T}{here}\mycalendar{Jan.'10}{05} will soon be something new to watch on the living room
TV: your relatives and friends in different parts of the world.

On Tuesday, Panasonic and LG Electronics, two of the top television makers, are to announce that
they are integrating the free online calling service Skype into their Internet-connected
high-definition televisions.

People who buy these TVs, along with an extra Web camera and microphone accessory designed for the
living room, can conduct free, live video chats and phone calls from the couch.

The announcement is the first of many expected from TV manufacturers at the Consumer Electronics
Show in Las Vegas this week. Television makers are trying to give consumers reasons to begin
replacing their high-definition TVs that they bought only in the last few years. TV makers are also
trotting out ever-slimmer TVs, adding Internet connections and preparing to introduce 3-D
technology.

``This is the year when Internet-connected TVs will start to take off, and there is no doubt that
soon every TV that ships will have built in Wi-Fi, webcams and microphones,'' said Jonathan
Rosenberg, chief technology strategist of Skype. ``This is our first step.''

The domestic TV market survived the recession in surprisingly good health, shipping 33.86 million
units during 2009, up 17 percent from 2008, while retail sales were generally flat, according to the
research firm iSuppli. But furious competition among manufacturers cut deeply into prices and
profits, with the price of a 46-inch LCD TV, for example, falling under \$1,000 for the first time.

Internet-enabled TV sets, though they have not yet proved a hit with consumers, can restore some of
those shrunken profits. A 50-inch plasma, high-definition TV with a broadband connection from LG
runs about \$300 more on Amazon.com than a TV of the same size without the connection. In some
cases, TV makers also get a slice of the revenue when customers make Web purchases from their TVs.

That is where Skype comes in. Until now, Web-connected TVs have accessed only a limited number of
online services, like widgets from Yahoo that offer weather and news updates, or Netflix's streaming
movie service. By adding other services and making a television more like a PC, TV makers now want
to change the very identity of the primary screen in the house.

``The TV is not just a one-way entertainment device, but a two-way communications device and a
portal into other people's lives,'' said Bob Perry, a senior vice president at Panasonic.

The Skype service on a TV will work much as it does on a PC, but with some limitations. A TV program
will stop playing once a Skype call is made or answered; TV processors are not yet powerful enough
to allow people to chat while they watch a show, the companies say.

Panasonic, based in Japan, and LG, based in South Korea, will sell specially designed Web cameras
for their Skype-enabled sets, which should cost \$100 to \$200. These cameras, unlike typical
webcams, are customized for the technically challenging environment of the living room, where there
can be a wide range of distances between the TV and viewers.

The deals give the newly independent Skype, formerly a division of eBay, another valuable foothold
in the home. About 520 million people around the world use the service to place free phone calls and
have video chats with one another from their computers and other devices that run Skype software.
The company makes money by charging competitive rates for people to call regular phone numbers and
for add-on services like voice mail.

Video chatting has become an increasingly important feature of Skype. The company says video chats
account for 34 percent of calls on the service, and as many as half on days like Christmas and New
Year's, when families seek to connect face to face.

To further nurture this activity, Skype is also announcing on Tuesday that its software for Windows
PCs and televisions will support high-definition video calling in the 720p HD format, if users have
webcams that support the technology.

Skype makes no money from video chats, at least for now. But Mr.~Rosenberg of the company says that
it may soon start charging for add-on services, like the ability to set up a video conference call
for three or more parties.

``What we finally created with the Skype network is a community of users that are available via
video. Having Skype on the TV is a big part of building that,'' he said.

Not everyone believes Internet-enabled TVs, with features like Skype, are going to make much of a
difference to consumers. Riddhi Patel, an analyst at iSuppli, says he thinks that innovations like
LED-backlighting for LCDs, which can result in a picture with a crisper contrast ratio and a thinner
screen, are what catch shoppers' eyes in stores.

``When a consumer walks into the store, they make a purchase when they can instantly see the benefit
of upgrading to a new TV,'' Ms.~Patel said.

Take Off Your Shoes, and Is the Parrot Loaded?

By JOE SHARKEY ``DOES he bite?'' the screener at the checkpoint asked warily.

``She doesn't bite,'' I said.

``Because we have to check under the wings,'' he said.

``In that case,'' I said, ``she might bite.''

At issue was our chatty little African Grey parrot, Rosie, who was watching the scene from inside
her travel cage at the security checkpoint at the Newark airport. This was last week, a few days
after a suspected terrorist tried to blow up an international flight on its descent into Detroit by
igniting some explosives hidden in his underwear.

While the explosion fizzled, it threw airport security into a tizzy.

My wife and I had never before flown with our two parrots, but this time they had to come along on
our nonstop flight to Phoenix. Rosie could fly in the cabin. But our other parrot, a blue-and-gold
macaw named Petey, is too big for the cabin and was already in the hands of the very helpful people
at Continental's PetSafe program, which transports pets in a heated, pressurized cargo section of
the plane. We wouldn't see him again for six hours.

We were very anxious at the checkpoint. My wife solved the problem, though. One of Rosie's tricks is
to spread her wings and lower her beak if you ask her to imitate an eagle.

``Rosie, do an eagle,'' my wife said. Inside her cage with the screener's face framed in the open
door, the bird promptly spread her wings wide.

The screener had his look under the wings and lowered his wand. Merriment ensued all around -- but
it had to look pretty silly.

And off we went to our flight to Arizona, which proceeded smoothly with the smaller parrot asleep in
her cage under a seat.

The final week of 2009, after the underwear bomber set fire to himself on that Delta-Northwest
flight approaching Detroit from Amsterdam, was one of the strangest periods yet in the annals of air
travel security since the 2001 terrorist attacks. Business travelers and other frequent fliers, who
had long adapted stoically to whatever new security hurdles were thrown in their paths, were no
longer reacting with equanimity.

In many cases, stoicism was replaced by ridicule over what some travelers saw as knee-jerk reactions
to the underwear bomber.

I heard from a bunch of business travelers. But my favorite e-mail message offered some
tongue-in-cheek security advice.

Mads Oyen, a policy specialist at Unicef in New York, suggested removing from a plane any specific
seat that had been used by a would-be terrorist. ``If he used, say, 36E, remove that seat. Then this
cannot be tried again,'' Mr.~Oyen wrote.

Another suggestion from Mr.~Oyen addressed the widely criticized Transportation Security
Administration rule, imposed right after the Christmas episode and quickly withdrawn, that
passengers on flights to the United States not be informed of where they were, by cockpit
announcements or by those video-screen maps that show a plane's position.

``Not only remove the in-flight maps, but tell people after seating that they are headed to a
destination they are not booked for,'' he wrote. ``It is important that security measures are
unpredictable.''

In a phone interview, Mr.~Oyen recalled that he once watched security screeners at an airport in
Uganda put a uniformed soldier's loaded gun through the X-ray scanner ``presumably to check for
sharp objects.''

On a more serious note, an airline pilot who did not want his name used, asked, ``When will
passengers say enough is enough with the ineffective theatrical security measures?''

That remains to be seen. As has already been reported, the T.S.A. again revised procedures on
Sunday, requiring pat-downs and extra screening for United States-bound passengers who are citizens
of 14 countries identified as being either state sponsors of terrorism or friendly toward certain
terrorists groups. At the same time, the T.S.A. relaxed the most strict security measures for United
States citizens and most others not flying through, or citizens of, those 14 nations.

Meanwhile, I should note that my wife and I and Rosie the parrot arrived on time in Phoenix, where
we then picked up Petey, the other parrot, at a special trailer set aside by Continental for
passengers to fetch pets.

When he heard my wife's voice, he let out a screech from his well-padded cage that could have been
heard in Las Vegas.

He had obviously had a good flight, too. And if he had any complaints about security, we didn't hear
about them. Petey wasn't talking.

\section{For F.D.R. Sleuths, New Focus on an Odd Spot}

\lettrine{W}{hen}\mycalendar{Jan.'10}{05} President Franklin D.~Roosevelt died unexpectedly on April
12, 1945, in Warm Springs, Ga., the White House lost no time announcing a cause of death.

The 63-year-old president, the shocked and grieving nation was told, had died of cerebral
hemorrhage. (``Last Words,'' read a front-page headline in The New York Times: `` 'I Have a Terrific
Headache.' '')

That Roosevelt died of a stroke is undisputed. But what caused it is a medical mystery that has
persisted to this day, a mystery heightened by the secrecy in which he, his aides and his doctors
always insisted on shrouding his health.

Now a new book -- ``F.D.R.'s Deadly Secret,'' by a neurologist, Dr.~Steven Lomazow, and a
journalist, Eric Fettmann (PublicAffairs) -- revives an intriguing theory.

Look closely at Roosevelt's portraits over his 12-year presidency. In his first two terms, there is
a dark spot over his left eyebrow. It seems to grow and then mysteriously vanishes sometime around
1940, leaving a small scar.

Was the spot a harmless mole? Or a cancerous melanoma that spread to contribute to, or even cause,
his death? Melanomas, after all, are known for causing strokes from bleeding when they spread to the
brain.

This hypothesis is not new. In 1979, Dr.~Harry S.~Goldsmith, then a surgeon at Dartmouth, wrote a
widely publicized medical journal article focusing attention on the possibility that the spot was a
melanoma. (I wrote an article about it at the time.) In 2007, after more medical sleuthing,
Dr.~Goldsmith published a book, ``A Conspiracy of Silence'' (iUniverse), fleshing out the theory.

What is different in the new book is the categorical claim that the killer was melanoma that
``metastasized to his brain, causing the growing tumor that would take Roosevelt's life a mere six
weeks later.''

But no matter how confidently the authors may assert it, the claim is still speculation -- unproved
and far from convincing.

Roosevelt's death was shocking in part because the White House and his doctors had kept secret how
sick he was. For example, though it was widely known that he had developed polio in 1921 at age 39,
he and his aides disguised the fact that he could not walk unaided and used heavy metal braces to
stand on paralyzed, withered legs. He used a wheelchair and demanded that photographers not show his
disabilities.

His terminal illness came during wartime, and in an era when leaders' health and other personal
matters were considered strictly private.

With rare exceptions, journalists were complicit. They did not probe the obvious clues they saw as
the president's appearance deteriorated.

Over his last year, for instance, he lost about 30 pounds. His doctors attributed it to a poor
appetite from a prescribed diet; Dr.~Lomazow and Mr.~Fettmann contend he became scrawny from a
spreading melanoma.

The authors point out that Turner Catledge, then a Washington correspondent for The New York Times
and later its executive editor, did not report how awful Roosevelt looked during an interview at the
White House in 1944, months before his nomination to an unprecedented fourth term.

Roosevelt was gaunt and glassy-eyed, Catledge wrote many years later; his jaw drooped, and he lost
his train of thought. Others witnessed similar episodes; in an interview, Dr.~Lomazow attributed
them to a type of seizure often associated with strokes.

Roosevelt's cardiologist, Dr.~Howard G.~Bruenn, certified that he died of a cerebral hemorrhage from
longstanding arteriosclerosis. Only in 1970 did Bruenn disclose in a medical journal article that
for many years the president's blood pressure was dangerously high. Available records show that it
had risen to 230/126 in 1944, from 128/82 in 1930, which would have contributed to heart failure. A
reading moments before he died was 300/190.

Even then, doctors knew that chronic high blood pressure (hypertension) and arteriosclerosis were a
potentially lethal combination that could cause heart disease and strokes. That became the standard
and most plausible explanation for Roosevelt's stroke.

The speculation about a melanoma cannot be verified because there was no autopsy and no known
biopsy, and most of Roosevelt's medical records disappeared shortly after his death from a safe in
the United States Naval Hospital in Bethesda, Md.

In their public accounts and the few surviving medical records, his doctors never suggested that
they performed a biopsy to determine whether he had any form of cancer. (Even during his lifetime
there were rumors that he had prostate cancer.)

Franklin Roosevelt was hardly the only president to bend the truth, if not lie, to hide his
illnesses from the public. Thirteen years before Roosevelt took office, President Woodrow Wilson, in
his second term, had a paralyzing stroke; Wilson's wife and aides kept that fact hidden from the
public while they took over the running of the government.

Fortunately, recent decades have yielded vast changes in medical practice and in perceptions of the
public's right to know about a political leader's health. But Roosevelt's ailments must still be
viewed in the context of the times.

In the 1930s and '40s, doctors had none of the now-standard blood pressure drugs like strong
diuretics, beta blockers and ACE inhibitors. They paid far less attention than they do now to moles
suspected of being melanomas; even if they had such suspicions, it is inconceivable that the spot
over Roosevelt's eye would have drawn the kind of attention focused during the 2008 presidential
campaign on the extensive surgery for the most serious of Senator John McCain's four melanomas.

Why did Roosevelt's spot vanish after 1940? Melanomas rarely regress on their own. But did his? Was
the spot removed for a biopsy? Or for cosmetic reasons? Despite the assertions in ``F.D.R.'s Deadly
Secret,'' the answers are unknown.

Even in the 1940s, some Washington insiders claimed that Roosevelt knew he was seriously ill, if not
dying, when he ran for his final term, and that was one of the reasons he dumped his vice president,
Henry A.~Wallace, from the Democratic ticket in favor of Harry S.~Truman. But most historians say he
picked Truman for political reasons, not for his qualifications.

In July 1944, Dr.~Frank H.~Lahey, a nationally prominent surgeon in Boston, consulted in Roosevelt's
case. In a memorandum to the record that was made public largely through Dr.~Goldsmith's efforts,
Dr.~Lahey said he told Roosevelt's White House physician, Adm. Ross T.~McIntire, that he doubted
Roosevelt's capacity to survive another term. But the memorandum did not mention cancer: it focused
on the president's failing heart.

Dr.~Lomazow and Mr.~Fettmann contend that McIntire would certainly have told the president he had
cancer. This is not clear; in that era doctors often withheld the word ``cancer'' from patients with
any form of the disease.

The authors build on a 2007 paper by Dr.~Barron H.~Lerner that challenged a 1970 account by Bruenn,
Roosevelt's cardiologist, of the president's illness and death. Like Dr.~Lerner, Dr.~Lomazow
contends that the account, which has been considered definitive, was incomplete.

One reason is that Bruenn did not mention problems like the vanished eyebrow spot and the blood
transfusions that Roosevelt needed to help correct a severe anemia in 1941. McIntire attributed the
anemia to bleeding hemorrhoids; Dr.~Lomazow contends it was something more serious.

After McIntire's death in 1959, Dr.~Lomazow said, ``it fell upon'' Bruenn to protect Roosevelt's
wishes to keep his health problems secret. A British physician, Dr.~Hugh L'Etang, was about to
publish a paper suggesting that Roosevelt might have had melanoma, Dr.~Lomazow said.

Also, he said, the Roosevelt family wanted Bruenn's cooperation in documenting that the president
had been mentally capable during the Allies' end-of-war conference at Yalta in February 1945. During
the cold war, detractors had taken to calling him ``the sick man at Yalta'' and saying Stalin had
taken advantage of him.

``F.D.R.'s Deadly Secret'' adds to the many accounts of how the president deceived the public about
his health. Challenging the conventional wisdom is admirable. But adopting an alternative theory as
fact requires convincing evidence that is lacking in this book.

In some places the book is inexcusably confusing. In making a point in one place, the authors say it
is speculative; but the same point in another sentence is stated as fact.

For example, journalists reported that Roosevelt departed a number of times from the prepared text
when he spoke to Congress after returning from Yalta. The authors of ``F.D.R.'s Deadly Secret''
examined short film clips and deduced that he had a defect in his left visual fields known as a
hemianopia. Although the authors say it is usually from a stroke, they state flatly that the visual
deficit was ``caused by a metastatic brain tumor.''

The authors say that though it is unclear whether Roosevelt's doctors fully understood the nature of
this postulated deficit, ``they certainly knew that the president's lesion was malignant and had
metastasized.'' The book says the abdominal pains Roosevelt experienced in his last year were
``caused by the cancer that had metastasized to his bowel.'' What is the proof?

In an interview, Dr.~Lomazow acknowledged the inconsistencies.

``That is a fair criticism of the way the book is written,'' he said. ``The book goes back and
forth, and if it led you to believe that this is incontrovertible absolute evidence, then you are
correct and I sincerely apologize for the misunderstanding.''

Some of the confusion, he said, may stem from the authors' and the publisher's failure to
cross-check what they wrote.

Dr.~Goldsmith, the author of the 1979 article raising the possibility that Roosevelt had a melanoma,
now says that after further research he doubts that was what killed him.

Roosevelt could have died from a stroke from his high blood pressure and also had a melanoma and
prostate cancer. But without additional records, it is impossible to be definitive about the cause
of death.

Regardless, his death and its aftermath make one thing clear: All presidents and their doctors
should make full disclosures about their health. As long as crucial facts are kept secret, theories,
conspiracies and hype may tarnish their image and long outlive them.

\section{Study Says Women With Mate Get Heavier}

\lettrine{I}{t}\mycalendar{Jan.'10}{05} is widely known that women tend to gain weight after giving
birth, but now a large study has found evidence that even among childless women, those who live with
a mate put on more pounds than those who live without one.

The differences, the scientists found, were stark.

After adjusting for other variables, the 10-year weight gain for an average 140-pound woman was 20
pounds if she had a baby and a partner, 15 if she had a partner but no baby, and only 11 pounds if
she was childless with no partner. The number of women with a baby but no partner was too small to
draw statistically significant conclusions.

There is no reason to believe that having a partner causes metabolic changes, so the weight gain
among childless women with partners was almost surely caused by altered behavior. Moreover, there
was a steady weight gain among all women over the 10 years of the study.

This does not explain the still larger weight gain in women who became pregnant. The lead author,
Annette J.~Dobson, a professor of biostatistics at the University of Queensland in Australia,
suggested that physiological changes might be at work.

``Women's bodies may adjust to the increased weight associated with having a baby,'' Dr.~Dobson
said. ``There may be a metabolic adjustment that goes on when women are pregnant that is hard to
reverse. This would be more consistent with our findings than any other explanation.''

The study covered more than 6,000 Australian women over a 10-year period ending in 2006.

At the start, the women ranged in age from 18 to 23. Each woman periodically completed a survey with
more than 300 questions about weight and height, age, level of education, physical activity, smoking
status, alcohol consumption, medications used and a wide range of other health and health care
issues.

By the end of the study, published in the January issue of The American Journal of Preventive
Medicine, more than half the women had college degrees, about three-quarters had partners and half
had had at least one baby. Almost all of the weight gain happened with the first baby; subsequent
births had little effect.

Also by the end of the study period, there were fewer smokers and risky drinkers than at the
beginning, more women who exercised less and a larger proportion without paid employment.

But even after adjusting for all of these factors and more, the differences in weight gain among
women with and without babies, and among women with and without partners, remained.

Despite the study's limitations -- weight was self-reported, for example, and the sample size
diminished over time because people dropped out -- other experts found the results valuable.

``It's interesting and brings out some important points,'' said Maureen A.~Murtaugh, an associate
professor of epidemiology at the University of Utah who has published widely on weight gain in
women. Perhaps, she suggested, a more active social life may help explain why women with partners
gain more weight.

``Think of going to a restaurant,'' Dr.~Murtaugh said. ``They serve a 6-foot man the same amount as
they serve me, even though I'm 5 feet 5 inches and 60 pounds lighter.''

The study included only women, but the researchers cited one earlier study that showed an increase
in obesity among men who had children, adding further evidence that social and behavioral factors
are part of the explanation.

Dr.~Dobson said the finding of weight gain among all the women, with families or without, was
troubling.

``This is a general health concern,'' she said. ``Getting married or moving in with a partner and
having a baby are events that trigger even further weight gain.

``From a prevention point of view, one can look at these as particular times when women need to be
especially careful.''

\section{Hong Kong Prepares to Accept Chinese Accounting Standards}

\lettrine{H}{ong}\mycalendar{Jan.'10}{05} Kong acceptance of Chinese accounting standards would be
an important advance in Beijing's drive to globalize its financial sector, but it could also
challenge international investors with reports prepared by an industry prone to scandal.

Under a proposed rule change likely to take effect next year, the Hong Kong stock exchange would let
locally listed Chinese companies file reports using their home accounting standards. The proposal is
designed to reduce costs and keep Hong Kong competitive with Shanghai.

But concerns about supervision of Chinese auditors have led to delays, making it unlikely that the
exchange will meet its Jan.~1 target date for implementing the change.

The change would also lead to different results in Hong Kong and China for certain industries, like
insurance, although those differences would be expected to fade over time as standards converged.

``There is a confidence issue,'' said Judy Wong, president of the Association of Chartered Certified
Accountants Hong Kong, a group whose members stand to lose substantial business to lower-cost
Chinese accounting firms under the rule change.

The Chinese Ministry of Finance ``should let the public know how they vet the applications from
mainland accounting and audit firms and what are the criteria,'' she said, adding that the
accountants' association agreed with the proposal's direction but contended that issues needed to be
resolved.

The unification of standards would let Chinese companies whose shares trade in Hong Kong, known in
the city as H-shares, post a single set of results for each reporting period identical to the
reports they put out for their China-listed shares.

Most Hong Kong-listed Chinese companies are expected to use their home standards for their Hong Kong
reports once the move has become official, both as a money-saving measure and to reduce confusion
that often occurs because of differences in the reports released in China and Hong Kong.

Insurance companies like China Life and Ping An Insurance now see some of the greatest variations
because of different treatment of investment gains and losses. Others, like banks, airlines and
resource companies, would see little or no change.

Investors do not expect any big effect on stock prices from the accounting change.

Patrick Yiu, a fund manager at Cash Asset Management in Hong Kong, said that he would not change his
investment strategy after the new accounting standards had been adopted and that the move should be
welcomed by the broader market.

``Though there is some concern over discipline among auditors in the mainland, I don't think it will
be a major issue, and many major Chinese enterprises will do their best to disclose the most useful
information to their investors,'' he said.

The Hong Kong proposal testifies to the growing importance of China for international investors, who
are increasingly looking at a steady Hong Kong that is likely to top global fund-raising through
initial public offerings this year at \$30 billion, thanks to a strong influx of new China listings.
It could raise another \$48 billion next year, even as Shanghai takes over the top spot with \$56
billion in new offerings, according to forecasts by Ernst \& Young.

Mainland accounting firms that prepare reports under Chinese standards charge about one-third what
their Hong Kong counterparts do, making them more attractive in terms of price.

But any cost gains could easily be offset by an industry where book-cooking scandals were so common
less than a decade ago that they led a former prime minister, Zhu Rongji, to call fraudulent
accounting a ``malignant tumor'' that threatened the country's economy.

To help minimize that risk, under the proposal it is now considering, Hong Kong would allow in only
Chinese auditors who have passed muster with China's Finance Ministry and the securities regulator,
the China Securities Regulatory Commission.

Other checks and balances are also being discussed.

``We expect it will come up with a list; it should not be more than five firms,'' said Jack Chow, a
partner based in Hong Kong at the global accounting firm KPMG.

That list is expected to include domestic leaders like RSM China, Shinewing Certified Public
Accountants and BDO China Shu Lun Pan C.P.A.'s, with future approvals expected to follow.

Hong Kong's move is also seen as an acknowledgement that Chinese accounting standards are rapidly
converging with international standards.

One major difference that remains between Chinese and international standards is accounting for
asset impairment, which is responsible for some of the big differences in results for
investment-oriented firms like insurers.

China's standards do not allow for the reversal of asset impairment if initial provisions later turn
out to have been overstated. But international accounting does allow for such reversal.

The Chinese are now working with their international counterparts to resolve that difference, said
Chris Joy, executive director of Hong Kong Institute of Certified Public Accountants.

``If convergence doesn't continue and if standards start to diverge, effectively the stock
exchange's proposal will fall aside,'' he said.

\section{China Starts to Lift Region's Web Blackout}

\lettrine{N}{early}\mycalendar{Jan.'10}{05} six months after gruesome ethnic rioting on the western
frontier of Xinjiang, Chinese authorities gingerly began to lift an Internet blackout on Tuesday,
allowing partial access to a pair of official news sites.

The Xinjiang regional government said it also planned to restore other Web sites and services,
cellphone text messaging and international telephone connections ``later, based on the relevant
situation, step by step,'' according to an official statement.

For now, though, people in Xinjiang can visit the sites of the official government news agency,
Xinhua, and the Communist Party's main newspaper, People's Daily. Even on those reliably policed
sites, the region's Web users were still barred from engaging in e-mail, blogging or forums. It also
took noticeably longer to load pages than it had before the riots, one user said.

After languishing under a communications lockdown as protracted and geographically far-reaching as
any in China in the digital age, some residents were thrilled with even so modest an opening.

``People here have been feeling so helpless and bored, you wouldn't believe how excited some were
today,'' said a Ms.~Li, a journalist in Urumqi, the capital of Xinjiang, who did not want to
disclose her full name for fear of losing her job.

The restrictions had ``played an important role in safeguarding the stability throughout the region,
but at the same time they had brought a good deal of inconvenience to normal production and life,''
the Xinjiang government acknowledged.

The authorities there severed Internet, cellphone messaging and international telephone service the
day after long-simmering frictions erupted on July 5 in Urumqi. Nearly 200 people were killed and
more than 1,700 were wounded in the rioting, according to official tolls, most of them Han Chinese.

The violence began when members of the Muslim, Turkic-speaking Uighur minority took to the streets,
partly in reaction to a deadly showdown in southeast China days before.

Chinese authorities charged that Uighur separatist networks mobilized rioters via phones and popular
social networking sites, and have defended the blackouts on security grounds.

But in Xinjiang, residents widely assert that the government squeeze on information stoked further
tensions and unrest, such as during a spate of syringe stabbings in Urumqi in August. Deadly
protests ensued in early September, and the city's mayor and the regional police chief were promptly
dismissed.

In Xinjiang, local authorities, banks and phone service providers have been able to send text
messages, but private citizens still cannot. People can read news on a number of local
government-run media sites that were restored in August, yet most of those sites are blocked to
viewers outside the region.

State firms and some large companies have been able to apply to hook up to the Web outside Xinjiang,
but smaller businesses generally cannot. Even at popular online trading posts that have been
switched back on, Web users in Xinjiang can deal only with others in Xinjiang. Some have driven
hundreds of miles away to Gansu Province to conduct business.

In September, Xinjiang passed a broadly worded bill banning online speech that incites separatism
and upsets national unity and social stability, and ordered service providers to monitor their
systems for such provocations. The authorities have enlisted local Communist Youth League members to
act as online ``supervisors.''

But even on local news portals that have been operating for months, Web forums, blog and e-mail
services remain off limits.

On Tuesday, the same constraints seemed to apply to the two newly accessible sites, Xinhuanet.com
and People.com.cn, people in Xinjiang said.

The government statement quoted a regional official saying e-mail services would eventually be
restored on major sites under certain ``conditions,'' but did not specify them. The statement asked
people for their further understanding and support.

\section{China Finds Huge Fraud by Officials}

\lettrine{C}{hinese}\mycalendar{Jan.'10}{05} officials misused or embezzled about \$35 billion in
government money in the first 11 months of the year, according to a national audit released this
week.

The announcement is the latest indication of how widespread corruption has become among government
agencies and how difficult it will be for Beijing to root it out.

The National Audit Office, which carried out the examination, did not disclose the size of the
budgets reviewed this year. But the agency, which is based in Beijing, said that it surveyed nearly
100,000 government departments and state-owned companies, and that more than 1,000 officials were
facing prosecution or disciplinary action because of the audits.

Auditors said government officials engaged in everything from money laundering and issuing
fraudulent loans to cheating the government through the sale or purchase of state land or mining
rights.

``Criminals are now more intelligent, and covert,'' Liu Jiayi, the director of the National Audit
Office, was quoted as saying in the state-run news media.

Prime Minister Wen Jiabao hailed the work of the auditors on Tuesday and called on them to monitor
government projects and prevent waste.

But analysts say the Communist Party faces significant hurdles in trying to curtail corruption.
Every year Beijing announces new anticorruption drives, new laws and new policies aimed at dealing
with the problem. But every year the scale of fraud seems enormous, particularly in a country where
the average person earns less than \$50 a week.

In 2005, for instance, the National Audit Office reported finding about \$35 billion worth of
government funds misused or embezzled. That was the last year the office gave a national figure
covering its audits, according to its Web site.

Experts say the audits revealed one thing: many in government are finding ways to steal public
money.

``The huge crackdown reflects the seriousness of corruption in China's government,'' said Zhu Lijia,
a professor of public policy at the Chinese Academy of Governance in Beijing. ``Even the National
Audit Office should be supervised. In the past few years it was the N.A.O. that decided whether to
publish or hide some statistics.''

\section{Clinton Urges Global Response to Internet Attacks}

\lettrine{D}{eclaring}\mycalendar{Jan.'10}{25} that an attack on one nation's computer networks
``can be an attack on all,'' Secretary of State Hillary Rodham Clinton issued a warning on Thursday
that the United States would defend itself from cyberattacks, though she left unclear the means of
response.

In a sweeping, pointed address that dealt with the Internet as a force for both liberation and
repression, Mrs.~Clinton said: ``Those who disrupt the free flow of information in our society or
any other pose a threat to our economy, our government and our civil society. Countries or
individuals that engage in cyber-attacks should face consequences and international condemnation.''

Her speech was the first in which a senior American official had articulated a vision for making
Internet freedom a plank of American foreign policy. While the details remained sketchy, her remarks
could have far-reaching consequences, given the confrontation between Google and the Chinese
government over the company's assertion that its networks had been subject to a sophisticated attack
that originated in mainland China.

Mrs.~Clinton called for China to investigate Google's accusation and be open about its findings. She
said that the United States supported Google in publicly defying the Chinese government's
requirement that it censor the contents of its Chinese-language search engine.

``Censorship should not be in any way accepted by any company from anywhere,'' Mrs.~Clinton said.
``American companies need to take a principled stand. This needs to be part of our national brand.''

This month Google announced that it was ``no longer willing to continue censoring'' search results
for its Chinese users, pointing to breaches of Gmail accounts held by human rights activists in
China. Several other companies had also been targets of hacking, the company found. Google has
avoided placing direct blame on the government in Beijing, which has sought to describe the
situation as strictly a business dispute.

The Obama administration has been similarly cautious. Last week, a senior administration official
said the United States would issue a ``d\'emarche'' -- a diplomatic move often used to lodge a
protest -- against China in the coming days. An official said Thursday that the administration would
hold off to see whether the Chinese responded to Mrs.~Clinton's call for an explanation of the
Google allegations.

The administration's dealings with China are further complicated by the American debt held by the
Chinese government and issues like climate change, on which the United States is seeking its
cooperation. Though Mrs.~Clinton said the administration would air its differences with Beijing, she
said it would be in the context of a ``positive, cooperative, and comprehensive relationship'' -- a
clause added to her speech at the last minute.

Mrs.~Clinton also identified Saudi Arabia, Egypt, Tunisia, Vietnam and Uzbekistan as countries that
constrain Internet freedom or persecute those who use the Web to circulate unpopular ideas. She
pointed to an Egyptian blogger, Bassem Samir, who was in the audience at the Newseum in Washington
for Mrs.~Clinton's speech and had been imprisoned by Egyptian authorities.

Human rights groups applauded the speech, though some questioned how the United States would enforce
the warnings.

Tom Malinowski, the Washington advocacy director for Human Rights Watch, said the United States
should treat China's forced censorship as an unfair trade practice, which could be confronted
through the World Trade Organization or raised in future trade negotiations.

Still, Mr.~Malinowski said: ``I really thought this was groundbreaking. She showed no hesitation in
naming countries, including U.S.~allies, for suppressing speech on the Internet. She made a very
strong case for connecting Internet freedom to core American national security interests.''

As secretary of state, Mrs.~Clinton has elevated the role of the Internet and digital technology in
American diplomacy. She named Alec Ross, a technology entrepreneur who advised the campaign of
President Obama, as her senior adviser for innovation.

Mr.~Ross has assembled a team that is pursuing programs like a social network for young people in
Pakistan and a service that lets people in Mexico file electronic reports on drug-related activity.

Mrs.~Clinton announced a new \$15 million effort to help more young people, women and citizens
groups in other countries communicate on the Web. None of the proposals she mentioned focused on
China or Iran, and the financing is relatively modest.

For Cameran Ashraf, 29, an Iranian-American information technology worker who has helped Iranian
protesters circumvent government filtering of their messages, Mrs.~Clinton's tone was enough. ``I
didn't expect such strong, forceful language,'' he said. ``I was beyond pleased.''

\section{China Rebuffs Clinton on Internet Warning}

\lettrine{T}{he}\mycalendar{Jan.'10}{25} Chinese Foreign Ministry lashed out Friday against
criticism of China in a speech on Internet censorship made by Secretary of State Hillary Rodham
Clinton, calling on the United States government ``to respect the truth and to stop using the
so-called Internet freedom question to level baseless accusations.''

Ma Zhaoxu, a Foreign Ministry spokesman, said in a written statement posted Friday afternoon on the
ministry's Web site that the criticism leveled by Mrs.~Clinton on Thursday was ``harmful to
Sino-American relations.''

``The Chinese Internet is open,'' he said.

The statement by the Foreign Ministry, along with a scathing editorial in the English-language
edition of The Global Times, a populist, patriotic newspaper, signaled that China was ready to
wrestle politically with the United States in the debate over Internet censorship.

President Obama promised last year to start a more conciliatory era in United States-China
relations, pushing human rights issues to the background, but the new criticism of China's Internet
censorship and rising tensions over currency valuation and Taiwan arms sales indicate that animus
could flare in the months ahead.

Mrs.~Clinton's sweeping speech with its cold war undertones -- likening the information curtain to
the Iron Curtain -- criticized several countries by name, including China, for Internet censorship.
It was the first speech in which a top administration official offered a vision for making Internet
freedom an integral part of foreign policy.

The debate over Internet censorship was brought to the fore in China last week when Google announced
it might shut down its Chinese-language search engine, Google.cn, and curtail its other operations
in mainland China if Chinese officials did not back down from requiring Google to censor search
results.

Until now, the Chinese government had been trying to frame the dispute with Google as a commercial
matter, perhaps because officials want to avoid having the dispute become a referendum on Internet
censorship policies among Chinese liberals and foreign companies operating in China. On Thursday, He
Yafei, a vice foreign minister, had said the Google dispute should not be ``over-interpreted'' or
linked to the bilateral relationship with the United States, according to Xinhua, the official state
news agency.

But in the aftermath of Mrs.~Clinton's speech, that attitude could be changing. Mrs.~Clinton
pointedly said that ``a new information curtain is descending across much of the world'' and
identified China as one of a handful of countries that had stepped up Internet censorship in the
past year. (Starting in late 2008, the Chinese government shut down thousands of Web sites under the
pretext of an antipornography campaign.) She also praised American companies such as Google that are
``making the issue of Internet and information freedom a greater consideration in their business
decisions.''

The State Department had invited at least two prominent Chinese bloggers to travel to Washington for
Mrs.~Clinton's speech, and on Friday the United States Embassy here invited bloggers, mostly
liberals, to attend a briefing on Internet issues.

A White House spokesman, Bill Burton, said Friday that ``all we are looking for from China are some
answers.''

In its editorial, the English-language edition of The Global Times said Mrs.~Clinton ``had raised
the stakes in Washington's clash with Beijing over Internet freedom.''

The American demand for an unfettered Internet was a form of ``information imperialism,'' the
newspaper said, because less developed nations cannot possibly compete with Western countries in the
arena of information flow.

``The U.S.~campaign for uncensored and free flow of information on an unrestricted Internet is a
disguised attempt to impose its values on other cultures in the name of democracy,'' the newspaper
said, adding that the ``U.S.~government's ideological imposition is unacceptable and, for that
reason, will not be allowed to succeed.''

Articles on the Chinese-language Web site of The Global Times asserted that the United States
employs the Internet as a weapon to achieve worldwide hegemony.

One big question is whether ordinary Chinese will, to any large degree, accept China's arguments
justifying Internet censorship. Although urban, middle-class Chinese often support government
policies on sovereignty issues such as Tibet or Taiwan, they generally deride media censorship. That
feeling is especially pronounced among those who call themselves netizens. China has the most
Internet users of any country, some 384 million by official count, but also the most complex system
of Internet censorship, nicknamed the Great Firewall.

Except in the western region of Xinjiang, which is only starting to restore Internet access after
cutting service off entirely after ethnic riots in July, canny netizens across China use software to
get over the Great Firewall while chafing at the controls.

\section{Nimble Agencies Sneak News Out of North Korea}

\lettrine{F}{or}\mycalendar{Jan.'10}{25} a journalist who helped break one of the biggest stories
out of North Korea in the past year, Mun Seong-hwi keeps an extremely low profile. The name he
offers is an alias. He does not reveal what he did in North Korea before his defection in 2006,
aside from mention of a ``desk job,'' in order to protect relatives left behind.

He also maintains a wall of secrecy around his three ``underground stringers'' in North Korea, who
he says do not know he works for Daily NK, an Internet news service based in Seoul and reviled by
Pyongyang.

On Nov.~30, quoting Mr.~Mun's and other anonymous ``sources inside North Korea,'' Daily NK reported
that, starting that day, the North Korean government would radically devalue its currency, requiring
people to exchange their old bank notes for new at a rate of 100 to 1. Furthermore, there would be
limit on how much of this old money people could turn in for new.

The report, which made headlines around the world and was later confirmed by South Korean officials,
had far-reaching implications. It meant, among other things, that the North Korean government was
cracking down on the country's nascent free markets, wiping out much of the wealth private
entrepreneurs had accumulated by trading goods at a time when the Communist government's ration
system was failing to meet its people's basic needs.

``I take pride in my work,'' Mr.~Mun, a man in his early 40s with brooding eyes and a receding
hairline, said in an interview. ``I help the outside world see North Korea as it is.''

Daily NK is one of six news outlets that have emerged in recent years specializing in collecting
information from North Korea. These Web sites or newsletters hire North Korean defectors and
cultivate sources inside a country shrouded in a near-total news blackout.

While North Korea shutters itself from the outside -- it blocks the Internet, jams foreign radio
broadcasts and monitors international calls -- it releases propaganda-filled dispatches through the
government's mouthpiece, the Korean Central News Agency.

But, thanks to Daily NK and the other services, it is also possible now for outsiders to read a
dizzying array of ``heard-in-North Korea'' reports, many on topics off limits for public discussion
in the North, like the health of the country's leader, Kim Jong-il.

The reports are sketchy at best, covering small pockets of North Korea society. Many prove wrong,
contradict each other or remain unconfirmed. But they have also produced important scoops, like the
currency devaluation and a recent outbreak of swine flu in North Korea. The mainstream media in
South Korea now regularly quote these cottage-industry news services.

``Technology made this possible,'' said Sohn Kwang-joo, the chief editor of Daily NK. ``We
infiltrate the wall of North Korea with cellphones.''

Over the past decade, the North's border with China has grown more porous as famine drove many North
Koreans out in search of food and an increasing traffic in goods -- and information -- developed. A
new tribe of North Korean merchants negotiates smuggling deals with Chinese partners, using Chinese
cellphones that pick up signals inside the North Korean border.

These phones have become a main tool of communication for many of the 17,000 North Korean defectors
living in the South trying to re-establish contact with their families and friends in the North.

Mr.~Sohn, a former reporter with the mainstream daily newspaper Dong-A in Seoul, has South Korean
``correspondents'' near the China-North Korea border.

These volunteers, many of them pro-democracy advocates during their student years, secretly meet
North Koreans traveling across the border and recruit underground stringers. The volunteers use
business visas, or sometimes pretend to be students or tourists.

``It's dangerous work, and it takes one or two years to recruit one,'' Mr.~Sohn said.

In the past year, the quality of the information these news services provide has improved as they
have hired more North Korean intellectuals and former officials who defected to the South and still
have friends in elite circles in the North, said Ha Tae-keung, a former student activist who runs
Open Radio for North Korea and a Web site.

``These officials provide news because they feel uncertain about the future of their regime and want
to have a link with the outside world, or because of their friendship with the defectors working for
us, or because of money,'' said Mr.~Ha, who also goes by his English name, Young Howard.

All these news outlets pay their informants. Mr.~Ha pays a bonus for significant scoops. Daily NK
and Open Radio each have 15 staff members, some of them defectors, and receive U.S.~congressional
funding through the National Endowment for Democracy, as well as support from other public and
private sources.

Recently, they have been receiving tips from North Koreans about corrupt officials.

``The fact that news comes out through civic groups like ours means that North Korean society is
changing fast,'' said Pomnyun Sumin, a Buddhist monk and chairman of Good Friends, a relief group
based in Seoul whose newsletter broke the swine flu story last month.

Some informants have become so adept with technology that they send text-messages, audio files and
photos to Seoul by cell phone, said Kim Heung-gwang, a former North Korean computer scientist who
heads North Korea Intellectuals Solidarity, a group of defectors that runs a news Web site.

Bringing news out of North Korea is risky. Mr.~Kim said that one of his informants was stopped last
May while trying to smuggle out a video in a small camera hidden in a cosmetics bottle. She is
believed to have killed herself in police detention, he said.

``You wonder whether they should let their sources take such risks,'' Chang Yong-hoon, who covers
North Korea for the mainstream news agency Yonhap, said at a recent forum on the news services.
``They've produced as many erroneous reports as they have real scoops.''

Kang Chol-hwan, a former North Korean prison camp inmate who now writes for the mainstream daily
Chosun, said there are ``information brokers'' in North Korea who sell exaggerated and fake news to
outside media. Lee Chan-ho, a chief analyst at the South Korean government's Unification Ministry,
warned that the ``flood of raw, unconfirmed reports'' complicates the effort to understand the
North.

Mr.~Ha, of Open Radio for North Korea, conceded that point: ``Because our sources have never been
trained in journalism, exaggeration is a problem for us. Some demand more money for information. We
try to cross-check our reports as much as possible.''

Mr.~Mun of Daily NK sends his stringers 1,000 Chinese renminbi, or about \$150, every two or three
months. They call him once a week at a designated time. They find a place where they feel safe from
the North Korean police patrolling with equipment that detects cell phone users, and dial. In Seoul,
Mr.~Mun hears the ring tone, then calls back and talks for about half an hour. After the call, his
stringers hang up and hide their phones until the next call.

Recently, with so many developments in the North, they have been calling him at unscheduled times,
for instance, when he was in the subway.

``I have to rush off and call back quickly,'' Mr.~Mun said. ``If I don't within five minutes, he'll
turn off and I'll lose him.''

There is another strict rule.

``We don't know, and never ask, each other's real name,'' Mr.~Mun said. ``That's safer for them.
Their safety is my biggest concern.''

\vocabulary
\end{document}
